% How accurate LIGHT v0_4 is: v0_3 with the 5p rung, the optics, KrF* + F,
% ion-ion recombination's limbs and bands, the walls and the wave. No
% number here is typed by hand. They come from
%   overview_v0_4_fast.tex, overview_v0_4_slow.tex   make_overview_v0_4.py
%   system_map_v0_4_numbers.tex                      make_system_map_v0_4.py
%   overview_v0_3_fast.tex, overview_v0_3_slow.tex   make_overview_v0_3.py
%   overview_v0_2_fast.tex, overview_v0_2_slow.tex   make_overview_v0_2.py
%   energy_ladder_v0_4.tex                           make_energy_ladder.py
% Build from LIGHT/docs:
%   python3 make_system_map_v0_4.py && python3 make_overview_v0_4.py
%   python3 make_energy_ladder.py v0_4
%   latexmk -pdf accuracy_overview_v0_4.tex
\documentclass[11pt]{article}
\usepackage[a4paper,margin=2.2cm]{geometry}
\usepackage[T1]{fontenc}
\usepackage[scaled=0.95]{helvet}
\renewcommand{\familydefault}{\sfdefault}
\usepackage{amsmath}
\usepackage{sansmath}
\usepackage{booktabs,tabularx,array}
\usepackage{xcolor}
\usepackage{float}
\usepackage{enumitem}
\usepackage{needspace}
\usepackage{tikz}
\usepackage{graphicx}
\usetikzlibrary{arrows.meta,decorations.pathmorphing}
\setlist{nosep, leftmargin=1.2em}
\setlength{\parindent}{0pt}
\setlength{\parskip}{5pt}
\definecolor{ink2}{HTML}{52514E}

% LIGHT v0_2 accuracy numbers (fast). GENERATED by make_overview_v0_2.py
% from LIGHT v0_2 (0.1.2) and ML/. Do not edit. Rerun the script instead.
\providecommand\widef[2]{\expandafter\def\csname wi@#1\endcsname{#2}}
\widef{xu.50}{0.91}
\widef{xu.50.v0}{2.89}
\widef{xu.100}{0.97}
\widef{xu.100.v0}{2.41}
\widef{xu.50.meas}{$5.0{\times}10^{7}$}
\widef{xu.100.meas}{$6.2{\times}10^{7}$}
\widef{eq5.dev}{$1.7{\times}10^{-4}$}
\widef{gap}{120}
\widef{closure.max}{$10^{-12}$}
\widef{closure.runs}{9}
\widef{hot.p}{0.1}
\widef{hot.P}{500}
\widef{hot.tg}{3458}
\widef{sc.eden_kr}{2.3}
\widef{sc.eden_f2}{3.2}
\widef{sc.ksu_fig2}{4.2}
\widef{sc.ksu_fig1}{6.1}
\widef{sc.xu_figs5_7}{11.7}
\widef{sc.hakuta_fig3}{34.1}
\widef{sc.hakuta_fig2}{19.3}
\widef{sc.ys_ir300}{\ensuremath{-}0.5}
\widef{sc.ys_ratio_230_300}{\ensuremath{-}8.4}
\widef{sc.bs_tau_b}{\ensuremath{-}22.3}
\widef{ys.meas}{0.16}
\widef{ys.err}{0.02}
\widef{ysr.meas}{1.30}
\widef{ysr.err}{0.06}
\widef{sc.train.lo}{2.3}
\widef{sc.train.hi}{6.1}
\widef{sc.bs.meas}{9.0}
\widef{sc.bs.err}{0.6}
\widef{tah.g0.meas}{6.39}
\widef{tah.g0}{5.69}
\widef{tah.g0.x}{0.89}
\widef{tah.isat.meas}{1.4}
\widef{tah.isat}{1.55}
\widef{tah.isat.err}{0.1}
\widef{tah.isat.x}{1.11}
\widef{tah.eta.meas}{0.27}
\widef{tah.eta}{0.268}
\widef{tah.eta.x}{0.99}
\widef{tah.form}{0.300}
\widef{tah.tg}{446}
\widef{tah.pin}{330}
\widef{tah.torr}{665}
\widef{tah.f2torr}{2.8}
\widef{tah.seed}{$1.0{\times}10^{10}$}
\widef{tah.def.g0}{5.69}
\widef{tah.win.lo}{200}
\widef{tah.win.hi}{400}
\widef{tah.f2left}{29}
\widef{tah.bd.kr2_f.lo}{5.69}
\widef{tah.bd.kr2_f.hi}{6.23}
\widef{tah.bd.trimer_tg_pow.lo}{5.69}
\widef{tah.bd.trimer_tg_pow.hi}{2.84}
\widef{tah.bd.kr_pen.lo}{5.71}
\widef{tah.bd.kr_pen.hi}{5.69}
\widef{tah.bd.kr2f_f2.lo}{5.69}
\widef{tah.bd.kr2f_f2.hi}{5.89}
\widef{tah.bd.kr2f_f.lo}{5.69}
\widef{tah.bd.kr2f_f.hi}{5.86}
\widef{tah.bd.kr2_form.lo}{5.69}
\widef{tah.bd.kr2_form.hi}{5.31}
\widef{tah.bd.kr2_f2.lo}{5.69}
\widef{tah.bd.kr2_f2.hi}{5.74}
\widef{tah.bd.kr2_pen.lo}{5.69}
\widef{tah.bd.kr2_pen.hi}{5.69}
\widef{tah.bd.minor}{0.38}
\widef{tah.bd.widest}{the trimer's $(T_g/300)^n$}
\widef{tah.hiF.g0}{6.23}
\widef{tah.hiF.x}{0.97}
\widef{tah.tri.g0}{2.84}
\widef{tah.tri.x}{0.44}
\widef{tri.lo.pow}{0.5}
\widef{tri.klimek.hi}{425}
\widef{tah.kannari.g0}{6.10}
\widef{tah.kannari.x}{0.95}
\widef{tah.own.g0}{5.53}
\widef{tah.own.te}{3.8}
\widef{tah.own.ne}{$2.59{\times}10^{14}$}
\widef{tah.own.tg}{423}
\widef{tah.beam.g0}{5.69}
\widef{tah.beam.te}{1.6}
\widef{tah.beam.ne}{$6.33{\times}10^{14}$}
\widef{tah.beam.tg}{446}
\widef{tah.own.x}{0.87}
\widef{kr2.def}{4.67}
\widef{kr2.vuv}{20.0}
\widef{kr2.burn}{640}
\widef{kr2.F}{$2.12{\times}10^{17}$}
\widef{kr2.lo.k}{$1.0{\times}10^{-11}$}
\widef{kr2.lo}{6.52}
\widef{kr2.kann.k}{$3.0{\times}10^{-10}$}
\widef{kr2.kann}{9.21}
\widef{kr2.kann.ratio}{17}
\widef{kr2.hi.k}{$1.0{\times}10^{-9}$}
\widef{kr2.hi}{9.50}
\widef{kr2.hi.vuv}{0.5}
\widef{kr2.hi.ratio}{57}
\widef{kr2.hi.afterpc}{74}
\widef{kr2.both}{9.17}
\widef{kr2.pen.k}{$5.0{\times}10^{-10}$}
\widef{pure.Kr.x}{1.02}
\widef{pure.Ar.x}{1.03}
\widef{pure.dr.lo}{1.00}
\widef{pure.dr.hi}{1.00}
\widef{pure.dr.te.lo}{0.026}
\widef{pure.dr.te.hi}{2}
\widef{pure.k15}{$2.4{\times}10^{-31}$}
\widef{pure.conv.lo}{$2.74{\times}10^{-31}$}
\widef{pure.conv.hi}{$3.3{\times}10^{-31}$}
\widef{pure.vuv}{0.57}
\widef{pure.vuv.meas}{0.42}
\widef{pure.vuv.err}{0.05}
\widef{pure.vuv.x}{1.36}
\widef{pure.vuv.ideal}{0.54}
\widef{pure.vuv.nothermal}{0.55}
\widef{pure.thermal}{0.060}
\widef{pure.band.move}{0.001}
\widef{pure.hi.P}{300}
\widef{pure.hi.pulse}{1}
\widef{pure.hi.torr}{665}
\widef{pure.hi.def}{0.27}
\widef{pure.hi.lo}{0.13}
\widef{pure.hi.lo.what}{\Krts{} + \Krts{} \ra \Krtp{} + 2Kr + \el{} at $5.0{\times}10^{-10}$\unit{cm^3\,s^{-1}}}
\widef{pure.hi.hi}{0.49}
\widef{pure.hi.hi.what}{\Krs{} + 2Kr \ra \Krts{} + Kr at $5.5{\times}10^{-32}$\unit{cm^6\,s^{-1}}}
\def\wiTahBands{\Krts{} + F \ra KrF(B) + Kr & 5.69 to 6.23 & \x{0.89} to \x{0.97}\\
the trimer's $(T_g/300)^n$ & 5.69 to 2.84 & \x{0.89} to \x{0.44}\\
\Krs{} + \Krs{} \ra \Krp{} + Kr + \el{} & 5.71 to 5.69 & \x{0.89} to \x{0.89}\\
\KtF{} + \Ft{} & 5.69 to 5.89 & \x{0.89} to \x{0.92}\\
\KtF{} + F \ra 2Kr + 2F & 5.69 to 5.86 & \x{0.89} to \x{0.92}\\
\Krs{} + 2Kr \ra \Krts{} + Kr, 300\unit{K} & 5.69 to 5.31 & \x{0.89} to \x{0.83}\\
\Krts{} + \Ft{}, both channels & 5.69 to 5.74 & \x{0.89} to \x{0.90}\\
\Krts{} + \Krts{} \ra \Krtp{} + 2Kr + \el{} & 5.69 to 5.69 & \x{0.89} to \x{0.89}\\}

% LIGHT v0_2 accuracy numbers (slow). GENERATED by make_overview_v0_2.py
% from LIGHT v0_2 (0.1.2) and ML/. Do not edit. Rerun the script instead.
\providecommand\widef[2]{\expandafter\def\csname wi@#1\endcsname{#2}}
\widef{fit.k_BC.lo}{$4.63{\times}10^{-12}$}
\widef{fit.k_BC.mid}{$9.8{\times}10^{-12}$}
\widef{fit.k_BC.hi}{$1.73{\times}10^{-11}$}
\widef{fit.K_eq.lo}{0.84}
\widef{fit.K_eq.mid}{1.07}
\widef{fit.K_eq.hi}{1.38}
\widef{fit.tau_C.lo}{41}
\widef{fit.tau_C.mid}{50}
\widef{fit.tau_C.hi}{63}
\widef{fit.tau_B.lo}{6.79}
\widef{fit.tau_B.mid}{7.02}
\widef{fit.tau_B.hi}{7.22}
\widef{fit.k_F2.lo}{$3.92{\times}10^{-10}$}
\widef{fit.k_F2.mid}{$4.57{\times}10^{-10}$}
\widef{fit.k_F2.hi}{$5.3{\times}10^{-10}$}
\widef{fit.k_F2_C.lo}{$2.74{\times}10^{-10}$}
\widef{fit.k_F2_C.mid}{$3.02{\times}10^{-10}$}
\widef{fit.k_F2_C.hi}{$3.32{\times}10^{-10}$}
\widef{fit.k3_Kr.lo}{$1.26{\times}10^{-30}$}
\widef{fit.k3_Kr.mid}{$1.92{\times}10^{-30}$}
\widef{fit.k3_Kr.hi}{$2.44{\times}10^{-30}$}
\widef{fit.k3_Kr_C.lo}{$1.93{\times}10^{-31}$}
\widef{fit.k3_Kr_C.mid}{$3.64{\times}10^{-31}$}
\widef{fit.k3_Kr_C.hi}{$5.12{\times}10^{-31}$}
\widef{fit.bshare.lo}{0.57}
\widef{fit.bshare.mid}{0.69}
\widef{fit.bshare.hi}{0.80}
\widef{fit.bshare.light}{0.68}
\widef{fit.n}{10}
\widef{post.g0.lo}{4.99}
\widef{post.g0.mid}{5.39}
\widef{post.g0.hi}{6.18}
\widef{post.n}{40}
\widef{link.base}{0.268}
\widef{link.yield.lo}{0.240}
\widef{link.yield.lo.d}{\ensuremath{-}0.028}
\widef{link.yield.hi}{0.282}
\widef{link.yield.hi.d}{+0.014}
\widef{link.yield.arf}{0.211}
\widef{link.yield.arf.d}{\ensuremath{-}0.057}
\widef{link.nascentC}{0.290}
\widef{link.nascentC.d}{+0.022}
\widef{link.ionY}{0.225}
\widef{link.ionY.d}{\ensuremath{-}0.043}
\widef{link.ionC}{0.251}
\widef{link.ionC.d}{\ensuremath{-}0.017}
\widef{link.ionRate}{0.266}
\widef{link.ionRate.d}{\ensuremath{-}0.002}
\widef{link.kr2F.kann}{0.285}
\widef{link.kr2F.kann.d}{+0.017}
\widef{link.kr2F.hi}{0.290}
\widef{link.kr2F.hi.d}{+0.022}
\widef{link.kr2split.lo}{0.260}
\widef{link.kr2split.lo.d}{\ensuremath{-}0.008}
\widef{link.kr2split.hi}{0.276}
\widef{link.kr2split.hi.d}{+0.008}
\widef{link.k3}{0.256}
\widef{link.k3.d}{\ensuremath{-}0.012}
\widef{link.minor}{0.013}
\widef{el.loop.te}{1.61}
\widef{el.loop.ne}{$6.49{\times}10^{14}$}
\widef{el.loop.g0}{5.63}
\widef{el.loop.form}{0.299}
\widef{el.light.te}{1.60}
\widef{el.light.ne}{$6.33{\times}10^{14}$}
\widef{el.light.g0}{5.69}
\widef{el.light.form}{0.300}
\widef{el.dev.te}{0.7}
\widef{el.dev.ne}{2.5}
\widef{el.dev.g0}{1.0}
\widef{el.dev.formation}{0.5}
\widef{el.eatt}{0.55}
\widef{el.eatt.32}{2.41}
\widef{el.rpow}{0.55}
\widef{el.rcap}{0.95}
\widef{sw.krw.lo}{0.96}
\widef{sw.krw.hi}{1.03}
\widef{sw.f2.lo}{0.48}
\widef{sw.f2.hi}{2.41}
\def\wiLinks{KrF$^*$ yield of \Krs{} + \Ft & 0.90 & 0.70 & the measured 0.90 less its stated 0.20 & 0.240 & \ensuremath{-}0.028\\
 &  & 1.0 & the published assignment & 0.282 & +0.014\\
 &  & 0.5 & as for Ar$^*$ + F$_2$ & 0.211 & \ensuremath{-}0.057\\
share of that KrF formed in C & 0.38 & 0 & if nascent C relaxes to B at 1 atm & 0.290 & +0.022\\
KrF$^*$ yield of ion-ion recombination & 1 & 0.64 & as for Cl$^-$ & 0.225 & \ensuremath{-}0.043\\
share of that KrF formed in C & 0 & 0.38 & as for SF$_6^-$ & 0.251 & \ensuremath{-}0.017\\
ion-ion recombination rate & Flannery & \x{0.38} & Avco's model input & 0.266 & \ensuremath{-}0.002\\
\Krts{} + F \ra KrF(B) + Kr & 0 & $3.0{\times}10^{-10}$ & Kannari's estimate & 0.285 & +0.017\\
 &  & $1.0{\times}10^{-9}$ & the top of its band & 0.290 & +0.022\\}

% LIGHT v0_3 accuracy numbers (fast). GENERATED by make_overview_v0_3.py
% from LIGHT v0_3 (0.1.3) and ML/. Do not edit. Rerun the script instead.
\providecommand\wkdef[2]{\expandafter\def\csname wk@#1\endcsname{#2}}
\wkdef{sc.eden_kr}{2.3}
\wkdef{sc.eden_f2}{3.2}
\wkdef{sc.ksu_fig2}{4.2}
\wkdef{sc.ksu_fig1}{6.1}
\wkdef{sc.xu_figs5_7}{11.7}
\wkdef{sc.hakuta_fig3}{34.1}
\wkdef{sc.hakuta_fig2}{19.3}
\wkdef{sc.ys_ir300}{\ensuremath{-}0.5}
\wkdef{sc.ys_ratio_230_300}{\ensuremath{-}8.4}
\wkdef{sc.bs_tau_b}{\ensuremath{-}22.3}
\wkdef{sc.same}{10}
\wkdef{sc.train.lo}{2.3}
\wkdef{sc.train.hi}{6.1}
\wkdef{xu.50}{0.91}
\wkdef{xu.100}{0.97}
\wkdef{ml.dF}{$2.2{\times}10^{-9}$}
\wkdef{ml.dk}{$4.2{\times}10^{-8}$}
\wkdef{ml.own}{$1.1{\times}10^{-5}$}
\wkdef{ml.states}{3}
\wkdef{bal.pulse}{$2.0{\times}10^{-6}$}
\wkdef{bal.pulse.nu}{0}
\wkdef{bal.nofield}{$5.6{\times}10^{-5}$}
\wkdef{bal.nofield.nu}{0}
\wkdef{bal.nogrowth}{$2.9{\times}10^{-4}$}
\wkdef{bal.nogrowth.raw}{$2.0{\times}10^{-3}$}
\wkdef{und.600.dev}{$1.4{\times}10^{-14}$}
\wkdef{und.600.cell}{0.029}
\wkdef{und.600.mean}{0.0377}
\wkdef{und.2400.dev}{$6.3{\times}10^{-14}$}
\wkdef{und.2400.cell}{0.007}
\wkdef{und.2400.mean}{0.0386}
\wkdef{und.32kt}{0.0388}
\wkdef{s.hot}{0.282}
\wkdef{s.cold}{0.042}
\wkdef{s.en.006}{0.444}
\wkdef{s.en.02}{0.388}
\wkdef{s.en.05}{0.473}
\wkdef{inj.dev}{$1.6{\times}10^{-5}$}
\wkdef{step.big}{$1.7{\times}10^{-12}$}
\wkdef{step.tiny}{$1.2{\times}10^{-10}$}
\wkdef{step.energy}{$2.8{\times}10^{-5}$}
\wkdef{step.ee600}{$1.0{\times}10^{-1}$}
\wkdef{step.ee2400}{$7.6{\times}10^{-3}$}
\wkdef{step.power}{$10^{9}$}
\wkdef{ref.p}{1}
\wkdef{ref.xf2}{0.5}
\wkdef{ref.P}{300}
\wkdef{ref.pulse}{1}
\wkdef{ref.tend}{2}
\wkdef{conv.dt.ns}{5}
\wkdef{und.lo}{600}
\wkdef{und.hi}{2400}
\wkdef{und.tg}{300}
\wkdef{s.t}{0.12}
\wkdef{s.field}{0.4}
\wkdef{s.m.006}{0.06}
\wkdef{s.m.02}{0.2}
\wkdef{s.m.05}{0.5}

% LIGHT v0_3 accuracy numbers (slow). GENERATED by make_overview_v0_3.py
% from LIGHT v0_3 (0.1.3) and ML/. Do not edit. Rerun the script instead.
\providecommand\wkdef[2]{\expandafter\def\csname wk@#1\endcsname{#2}}
\wkdef{ref.248}{4.79}
\wkdef{ref.248.v02}{4.67}
\wkdef{ref.248.d}{+2.7}
\wkdef{ref.272}{0.153}
\wkdef{ref.272.v02}{0.149}
\wkdef{ref.272.d}{+2.7}
\wkdef{ref.275}{0.735}
\wkdef{ref.275.v02}{0.712}
\wkdef{ref.275.d}{+3.2}
\wkdef{ref.420}{4.96}
\wkdef{ref.420.v02}{4.87}
\wkdef{ref.420.d}{+1.8}
\wkdef{ref.147}{19.7}
\wkdef{ref.147.v02}{20.0}
\wkdef{ref.147.d}{\ensuremath{-}1.8}
\wkdef{ref.all}{30.3}
\wkdef{ref.all.v02}{30.4}
\wkdef{ref.tg}{650}
\wkdef{ref.tg.v02}{649}
\wkdef{ref.krstar}{99.4}
\wkdef{ref.krstar.v02}{98.3}
\wkdef{ref.krplus}{26.0}
\wkdef{ref.krplus.v02}{26.8}
\wkdef{ref.krfb}{24.6}
\wkdef{ref.krfb.v02}{23.8}
\wkdef{ref.elastic}{5.3}
\wkdef{ref.elastic.v02}{6.2}
\wkdef{ref.closure}{$10^{-13}$}
\wkdef{ref.solutions}{355}
\wkdef{ref.en.lo}{0.98}
\wkdef{ref.en.hi}{43}
\wkdef{ref.ratio.lo}{1.006}
\wkdef{ref.ratio.hi}{1.031}
\wkdef{ref.gap}{$2.9{\times}10^{-5}$}
\wkdef{conv.step}{$3.9{\times}10^{-5}$}
\wkdef{conv.dt}{$6.8{\times}10^{-5}$}
\wkdef{conv.cells}{$1.6{\times}10^{-4}$}
\wkdef{ref.ratio.dt.hi}{1.015}
\wkdef{first.d.steady}{$1.9{\times}10^{-4}$}
\wkdef{first.d.maxwell}{$9.6{\times}10^{-4}$}
\wkdef{first.n}{87}
\wkdef{first.steady.n}{88}
\wkdef{first.steady.goff}{86}
\wkdef{first.mem.goff}{0}
\wkdef{first.mem.unreached}{0}
\wkdef{first.steady.ratio.lo}{0.02}
\wkdef{first.steady.ratio.hi}{0.19}
\wkdef{first.mem.ratio.lo}{0.98}
\wkdef{first.mem.ratio.hi}{1.07}
\wkdef{first.steady.gap}{0.22}
\wkdef{first.mem.gap0}{0.07}
\wkdef{first.mem.gap}{$1.5{\times}10^{-3}$}
\wkdef{first.edep}{$1.0{\times}10^{-3}$}
\wkdef{first.t}{1}
\wkdef{first.t.split}{0.1}
\wkdef{closure.max}{$10^{-13}$}
\wkdef{closure.runs}{4}
\wkdef{fac.ref.te}{3.00}
\wkdef{fac.ref.t}{502}
\wkdef{fac.tah.te}{1.61}
\wkdef{fac.tah.t}{302}
\wkdef{tah.g0}{5.64}
\wkdef{tah.g0.x}{0.88}
\wkdef{tah.g0.meas}{6.39}
\wkdef{tah.isat}{1.56}
\wkdef{tah.isat.x}{1.11}
\wkdef{tah.isat.meas}{1.4}
\wkdef{tah.eta}{0.266}
\wkdef{tah.eta.x}{0.99}
\wkdef{tah.eta.meas}{0.27}
\wkdef{tah.form}{0.299}
\wkdef{tah.win.g0}{5.64}
\wkdef{tah.win.te}{1.61}
\wkdef{tah.win.ne}{$6.52{\times}10^{14}$}
\wkdef{tah.win.tg}{446}
\wkdef{tah.gap}{$6.0{\times}10^{-4}$}
\wkdef{tah.reached}{every}
\wkdef{tah.en}{0.13}
\wkdef{first.tah.d}{$3.6{\times}10^{-4}$}
\wkdef{first.tah.ne.d}{$9.2{\times}10^{-4}$}
\wkdef{first.tah.steady.unreached}{80}
\wkdef{first.tah.mem.unreached}{0}
\wkdef{tah.kannari.g0}{6.05}
\wkdef{tah.kannari.x}{0.95}
\wkdef{tah.bd.kr2_f.lo}{5.64}
\wkdef{tah.bd.kr2_f.hi}{6.18}
\wkdef{tah.bd.trimer_tg_pow.lo}{5.64}
\wkdef{tah.bd.trimer_tg_pow.hi}{2.82}
\wkdef{tah.bd.kr_pen.lo}{5.66}
\wkdef{tah.bd.kr_pen.hi}{5.64}
\wkdef{tah.bd.kr2f_f2.lo}{5.64}
\wkdef{tah.bd.kr2f_f2.hi}{5.84}
\wkdef{tah.bd.kr2f_f.lo}{5.64}
\wkdef{tah.bd.kr2f_f.hi}{5.81}
\wkdef{tah.bd.kr2_form.lo}{5.64}
\wkdef{tah.bd.kr2_form.hi}{5.26}
\wkdef{tah.bd.kr2_f2.lo}{5.64}
\wkdef{tah.bd.kr2_f2.hi}{5.69}
\wkdef{tah.bd.kr2_pen.lo}{5.64}
\wkdef{tah.bd.kr2_pen.hi}{5.64}
\wkdef{tah.bd.widest}{the trimer's $(T_g/300)^n$}
\wkdef{tah.hiF.g0}{6.18}
\wkdef{tah.hiF.x}{0.97}
\wkdef{tah.tri.g0}{2.82}
\wkdef{tah.tri.x}{0.44}
\wkdef{el.loop.te}{1.610}
\wkdef{el.loop.ne}{$6.49{\times}10^{14}$}
\wkdef{el.loop.g0}{5.632}
\wkdef{el.loop.form}{0.2988}
\wkdef{el.own.te}{1.609}
\wkdef{el.own.ne}{$6.52{\times}10^{14}$}
\wkdef{el.own.g0}{5.637}
\wkdef{el.own.form}{0.2990}
\wkdef{el.dev.te}{\ensuremath{-}0.08}
\wkdef{el.abs.te}{0.08}
\wkdef{el.dev.ne}{+0.47}
\wkdef{el.abs.ne}{0.47}
\wkdef{el.dev.g0}{+0.08}
\wkdef{el.abs.g0}{0.08}
\wkdef{el.dev.formation}{+0.07}
\wkdef{el.abs.formation}{0.07}
\wkdef{el.passes}{10}
\wkdef{pure.vuv}{0.55}
\wkdef{pure.vuv.meas}{0.42}
\wkdef{pure.vuv.err}{0.05}
\wkdef{pure.vuv.x}{1.32}
\wkdef{pure.hi}{0.28}
\wkdef{pure.hi.P}{300}
\def\wkFactors{G1 & 0.076 & 0.070\\
C1a & 0.013 & $4.8{\times}10^{-3}$\\
C1b & $3.4{\times}10^{-3}$ & $7.0{\times}10^{-4}$\\
C1c & $1.5{\times}10^{-4}$ & $1.2{\times}10^{-5}$\\
C2 & $2.0{\times}10^{-5}$ & $2.3{\times}10^{-7}$\\
ATT & 0.58 & 0.96\\
DIS\_a3 & 1.2 & 1.0\\
DIS\_A1 & 1.1 & 0.97\\
VIB1 & 0.65 & 0.97\\
VIB4 & 0.69 & 0.98\\
R11p & 0.87 & 0.98\\}
\def\wkTahBands{\Krts{} + F \ra KrF(B) + Kr & 5.64 to 6.18 & \x{0.88} to \x{0.97}\\
the trimer's $(T_g/300)^n$ & 5.64 to 2.82 & \x{0.88} to \x{0.44}\\
\Krs{} + \Krs{} \ra \Krp{} + Kr + \el{} & 5.66 to 5.64 & \x{0.89} to \x{0.88}\\
\KtF{} + \Ft{} & 5.64 to 5.84 & \x{0.88} to \x{0.91}\\
\KtF{} + F \ra 2Kr + 2F & 5.64 to 5.81 & \x{0.88} to \x{0.91}\\
\Krs{} + 2Kr \ra \Krts{} + Kr, 300\unit{K} & 5.64 to 5.26 & \x{0.88} to \x{0.82}\\
\Krts{} + \Ft{}, both channels & 5.64 to 5.69 & \x{0.88} to \x{0.89}\\
\Krts{} + \Krts{} \ra \Krtp{} + 2Kr + \el{} & 5.64 to 5.64 & \x{0.88} to \x{0.88}\\}

% LIGHT v0_4 map numbers. GENERATED by make_system_map_v0_4.py
% from the model (0.1.4). Do not edit. Rerun the script instead.
\providecommand\wmdef[2]{\expandafter\def\csname wm@#1\endcsname{#2}}
\wmdef{e.P}{100.0}
\wmdef{e.pen}{1.2}
\wmdef{e.pen2}{0.0}
\wmdef{e.G1}{54.0}
\wmdef{e.C1}{6.7}
\wmdef{e.C2}{0.2}
\wmdef{e.F2}{2.7}
\wmdef{e.76}{33.8}
\wmdef{e.stepS}{0.0}
\wmdef{e.stepP}{0.0}
\wmdef{e.G6}{3.4}
\wmdef{e.dU}{0.0}
\wmdef{e.cap}{0.5}
\wmdef{C1.kr}{6.5}
\wmdef{C1.heat}{0.2}
\wmdef{F2e.F}{1.3}
\wmdef{F2e.heat}{1.8}
\wmdef{F2e.anion}{-0.4}
\wmdef{kpp.in}{290.0}
\wmdef{kpp.back}{4.2}
\wmdef{KPS.heat}{34.1}
\wmdef{PRAD.hv}{0.3}
\wmdef{PRAD.heat}{0.0}
\wmdef{kpp.harp}{3.1}
\wmdef{kpp.stepP}{0.0}
\wmdef{kpp.left}{0.0}
\wmdef{brp.KPS}{96.1}
\wmdef{brp.PRAD}{0.8}
\wmdef{brp.76r}{2.1}
\wmdef{brp.harp}{1.1}
\wmdef{harpP.B}{0.9}
\wmdef{harpP.C}{0.5}
\wmdef{harpP.F}{0.2}
\wmdef{harpP.heat}{1.4}
\wmdef{krs.in}{305.7}
\wmdef{krs.harp}{43.4}
\wmdef{krs.G2}{4.2}
\wmdef{krs.AK8}{10.5}
\wmdef{krs.stepS}{0.0}
\wmdef{krs.left}{0.0}
\wmdef{brs.harp}{14.2}
\wmdef{brs.G2}{1.4}
\wmdef{brs.AK8}{3.4}
\wmdef{brs.76}{81.0}
\wmdef{harp.B}{14.3}
\wmdef{harp.C}{8.8}
\wmdef{harp.F}{3.9}
\wmdef{harp.heat}{16.5}
\wmdef{G2.ion}{3.0}
\wmdef{C2.ion}{0.2}
\wmdef{stepS.ion}{0.0}
\wmdef{stepP.ion}{0.0}
\wmdef{AK8.heat}{0.8}
\wmdef{k2s.in}{9.7}
\wmdef{k2s.RAD}{4.9}
\wmdef{k2s.F2}{4.7}
\wmdef{k2s.F}{0.0}
\wmdef{k2s.pen}{0.0}
\wmdef{k2s.left}{0.1}
\wmdef{brk2.RAD}{50.9}
\wmdef{brk2.F2}{48.6}
\wmdef{RAD.hv}{4.5}
\wmdef{RAD.heat}{0.4}
\wmdef{NK7.B}{1.5}
\wmdef{BR.k2f}{1.4}
\wmdef{K2F2.F}{0.4}
\wmdef{K2F2.heat}{1.5}
\wmdef{k2F.B}{0.0}
\wmdef{k2F.Fin}{0.0}
\wmdef{k2F.heat}{0.0}
\wmdef{pen2.ion}{0.0}
\wmdef{K15.heat}{0.3}
\wmdef{R11p.in}{1.5}
\wmdef{R11p.kr}{1.3}
\wmdef{R11p.heat}{0.2}
\wmdef{ii.in}{2.2}
\wmdef{ii.anion}{-0.4}
\wmdef{ii.B}{1.0}
\wmdef{ii.C}{0.0}
\wmdef{ii.X}{0.0}
\wmdef{ii.heat}{0.8}
\wmdef{ions.left}{-0.7}
\wmdef{anion.left}{0.0}
\wmdef{mix.net}{2.9}
\wmdef{mix.netC}{2.9}
\wmdef{mix.heat}{0.0}
\wmdef{mix.BC}{12.5}
\wmdef{mix.CB}{15.4}
\wmdef{B.G4}{6.8}
\wmdef{B.G4A}{0.2}
\wmdef{B.G5}{12.1}
\wmdef{B.Q}{1.4}
\wmdef{B.QF2}{1.0}
\wmdef{B.QKR}{0.0}
\wmdef{B.QE}{0.4}
\wmdef{B.QF}{0.0}
\wmdef{B.left}{0.0}
\wmdef{C.G4C}{1.1}
\wmdef{C.G5C}{3.0}
\wmdef{C.Q}{2.3}
\wmdef{C.QF2}{0.8}
\wmdef{C.QKR}{1.1}
\wmdef{C.QE}{0.4}
\wmdef{C.QF}{0.0}
\wmdef{C.left}{0.0}
\wmdef{G4.hv}{5.8}
\wmdef{G4.F}{0.9}
\wmdef{G4.heat}{0.1}
\wmdef{G4A.hv}{0.2}
\wmdef{G4A.F}{0.0}
\wmdef{G4A.heat}{0.0}
\wmdef{G4C.hv}{0.8}
\wmdef{G4C.F}{0.1}
\wmdef{G4C.heat}{0.1}
\wmdef{QB.F}{0.5}
\wmdef{QB.heat}{1.0}
\wmdef{QC.F}{0.5}
\wmdef{QC.heat}{1.8}
\wmdef{tri.k2f}{13.6}
\wmdef{tri.heat}{1.5}
\wmdef{k2f.C3}{6.9}
\wmdef{k2f.Q2F}{8.1}
\wmdef{k2f.Q2FF}{0.0}
\wmdef{k2f.left}{0.0}
\wmdef{brk.C3}{45.8}
\wmdef{brk.Q2F}{54.1}
\wmdef{C3.hv}{3.8}
\wmdef{C3.F}{1.0}
\wmdef{C3.heat}{2.0}
\wmdef{Q2F.F}{1.2}
\wmdef{Q2F.heat}{6.9}
\wmdef{Q2FF.F}{0.0}
\wmdef{Q2FF.heat}{0.0}
\wmdef{FF.in}{0.0}
\wmdef{FF.heat}{0.0}
\wmdef{out.248}{5.8}
\wmdef{out.272}{0.2}
\wmdef{out.275}{0.8}
\wmdef{out.420}{3.8}
\wmdef{out.147}{4.5}
\wmdef{out.IR}{0.3}
\wmdef{out.heat}{71.6}
\wmdef{out.elastic}{3.4}
\wmdef{out.ions}{-0.7}
\wmdef{out.F}{10.2}
\wmdef{out.other}{0.1}
\wmdef{G4x.F}{1.0}
\wmdef{G4x.heat}{0.1}
\wmdef{iiNK7.B}{2.5}
\wmdef{k2s.fates}{9.7}
\wmdef{KRx.heat}{1.9}
\wmdef{harpx.F}{4.1}
\wmdef{harpx.heat}{17.9}
\wmdef{harpx.B}{15.2}
\wmdef{harpx.C}{9.3}
\wmdef{version}{0.1.4}
\wmdef{p}{1}
\wmdef{xF2}{0.5}
\wmdef{Tg0}{300}
\wmdef{P}{300}
\wmdef{pulse}{1}
\wmdef{tend}{2}
\wmdef{ne}{10^{15}}
\wmdef{neF2}{1.13{\times}10^{14}}
\wmdef{nePeak}{1.00{\times}10^{15}}
\wmdef{N}{2.45{\times}10^{19}}
\wmdef{Edep}{0.30}
\wmdef{Te}{5.81}
\wmdef{TeMid}{5.76}
\wmdef{rTail}{0.10}
\wmdef{TtMid}{0.57}
\wmdef{TgEnd}{742}
\wmdef{TgMax}{742}
\wmdef{pEnd}{2.48}
\wmdef{F2burn}{848}
\wmdef{F2left}{2.55}
\wmdef{krsFormed}{5.77{\times}10^{17}}
\wmdef{kppFormed}{4.81{\times}10^{17}}
\wmdef{k2sFormed}{1.98{\times}10^{16}}
\wmdef{krfFormed}{8.63{\times}10^{16}}
\wmdef{k2fFormed}{5.32{\times}10^{16}}
\wmdef{krf.eff}{23.0}
\wmdef{frm.ii}{3.6}
\wmdef{frm.kr2}{5.6}
\wmdef{frm.5p}{5.3}
\wmdef{brf.G4}{25.1}
\wmdef{brf.G4A}{0.9}
\wmdef{brf.G4C}{4.0}
\wmdef{brf.tri}{56.1}
\wmdef{brf.e}{3.1}
\wmdef{brf.F2}{6.8}
\wmdef{brf.Kr}{4.1}
\wmdef{brf.F}{0.0}
\wmdef{brf.QB}{5.4}
\wmdef{brf.QC}{8.5}
\wmdef{ion.R11p}{41.9}
\wmdef{ion.ii}{58.1}
\wmdef{brs.AK8.F2}{3.3}
\wmdef{brs.76.F2}{77.7}
\wmdef{cyc.gross}{0.8}
\wmdef{cyc.e}{33.8}
\wmdef{k.harp}{8.47{\times}10^{-10}}
\wmdef{k.harpP}{1.3{\times}10^{-9}}
\wmdef{pc.harpY}{90}
\wmdef{pc.harpB}{56}
\wmdef{pc.harpC}{34}
\wmdef{pc.harpX}{10}
\wmdef{pc.fC}{38}
\wmdef{k.G2}{5.0{\times}10^{-10}}
\wmdef{tauB}{6.99}
\wmdef{pc.BA}{3.5}
\wmdef{tauC}{50}
\wmdef{tauC3}{220}
\wmdef{gapBC}{120}
\wmdef{k.BC}{1.05{\times}10^{-11}}
\wmdef{Keq}{1.10}
\wmdef{k.G5}{1.9{\times}10^{-30}}
\wmdef{k.G5C}{3.79{\times}10^{-31}}
\wmdef{tpow.G5}{-3.32}
\wmdef{k.Q2F}{1.2{\times}10^{-10}}
\wmdef{k.K15}{2.4{\times}10^{-31}}
\wmdef{tpow.K15}{-0.75}
\wmdef{k.R11p}{1.6{\times}10^{-6}}
\wmdef{pow.R11p}{0.55}
\wmdef{cap.R11p}{0.95}
\wmdef{k.FF}{2.8{\times}10^{-34}}
\wmdef{n3.ii}{-3}
\wmdef{k3.ii}{4.82{\times}10^{-6}}
\wmdef{kL.ii}{5.66{\times}10^{-6}}
\wmdef{k.AK8}{1.7{\times}10^{-32}}
\wmdef{tpow.AK8}{-0.64}
\wmdef{tauRAD}{270}
\wmdef{k.kr2F2}{2.1{\times}10^{-10}}
\wmdef{pc.NK7}{50}
\wmdef{E.kr2s}{9.20}
\wmdef{E.kr2s.bind}{0.715}
\wmdef{hv.RAD}{8.43}
\wmdef{E.PP}{11.30}
\wmdef{gap5p}{1.385}
\wmdef{k.KPS}{1.7{\times}10^{-10}}
\wmdef{tauPRAD}{29}
\wmdef{A.76}{9.0{\times}10^{-7}}
\wmdef{B.76}{0.50}
\wmdef{C.76}{1.51}
\wmdef{C.stepS}{4.04}
\wmdef{hv.PRAD}{1.385}
\wmdef{E.anion}{2.6}
\wmdef{hv.G4C}{4.51}
\wmdef{hv.G4}{5.00}
\wmdef{hv.G4A}{4.56}
\wmdef{hv.C3}{2.95}
\wmdef{E.B}{5.856}
\wmdef{notes.krs}{60.6}
\wmdef{notes.krp}{3.2}
\wmdef{notes.emit}{15.5}
\wmdef{nodes}{47}
\wmdef{op.t0}{500}
\wmdef{op.t1}{1000}
\wmdef{op.g0}{4.12}
\wmdef{op.alpha}{0.63}
\wmdef{op.net}{3.50}
\wmdef{op.g0peak}{5.10}
\wmdef{op.tpeak}{560}
\wmdef{op.by}{\KtF{} 0.44, \Krpp{} 0.09, \Krs{} 0.03, \Ft{} 0.02, \Krtp{} 0.02, \Fm{} 0.02}
\wmdef{op.sigB}{2.6{\times}10^{-16}}
\wmdef{op.sigK2F}{1.0{\times}10^{-18}}
\wmdef{om.electron}{5.1}
\wmdef{om.basic}{5.8}
\wmdef{om.v0}{16.5}
\wmdef{om.v03}{4.67}
\wmdef{om.kannari}{5.18}
\wmdef{bd.trimer_tg_pow.def}{-3.32}
\wmdef{bd.trimer_tg_pow.lo}{-3.32}
\wmdef{bd.trimer_tg_pow.hi}{-0.50}
\wmdef{bd.trimer_tg_pow.lo248}{5.79}
\wmdef{bd.trimer_tg_pow.hi248}{2.29}
\wmdef{bd.krf_f.def}{0}
\wmdef{bd.krf_f.lo}{0}
\wmdef{bd.krf_f.hi}{7.8{\times}10^{-10}}
\wmdef{bd.krf_f.lo248}{5.79}
\wmdef{bd.krf_f.hi248}{3.99}
\wmdef{bd.kr2_f.def}{0}
\wmdef{bd.kr2_f.lo}{0}
\wmdef{bd.kr2_f.hi}{1.0{\times}10^{-9}}
\wmdef{bd.kr2_f.lo248}{5.79}
\wmdef{bd.kr2_f.hi248}{7.59}
\wmdef{bd.k_5s5p.def}{9.0{\times}10^{-7}}
\wmdef{bd.k_5s5p.lo}{0}
\wmdef{bd.k_5s5p.hi}{9.0{\times}10^{-7}}
\wmdef{bd.k_5s5p.lo248}{4.67}
\wmdef{bd.k_5s5p.hi248}{5.79}
\wmdef{bd.kstep_s.def}{0}
\wmdef{bd.kstep_s.lo}{0}
\wmdef{bd.kstep_s.hi}{1.5{\times}10^{-7}}
\wmdef{bd.kstep_s.lo248}{5.79}
\wmdef{bd.kstep_s.hi248}{5.13}
\wmdef{bd.kr2_form.def}{1.7{\times}10^{-32}}
\wmdef{bd.kr2_form.lo}{1.7{\times}10^{-32}}
\wmdef{bd.kr2_form.hi}{5.5{\times}10^{-32}}
\wmdef{bd.kr2_form.lo248}{5.79}
\wmdef{bd.kr2_form.hi248}{5.31}
\wmdef{bd.kstep_p.def}{0}
\wmdef{bd.kstep_p.lo}{0}
\wmdef{bd.kstep_p.hi}{2.0{\times}10^{-7}}
\wmdef{bd.kstep_p.lo248}{5.79}
\wmdef{bd.kstep_p.hi248}{5.56}
\wmdef{bd.kr_pen.def}{5.0{\times}10^{-10}}
\wmdef{bd.kr_pen.lo}{1.5{\times}10^{-10}}
\wmdef{bd.kr_pen.hi}{5.0{\times}10^{-10}}
\wmdef{bd.kr_pen.lo248}{6.01}
\wmdef{bd.kr_pen.hi248}{5.79}
\wmdef{bd.kr2f_f2.def}{1.2{\times}10^{-10}}
\wmdef{bd.kr2f_f2.lo}{1.2{\times}10^{-10}}
\wmdef{bd.kr2f_f2.hi}{4.3{\times}10^{-10}}
\wmdef{bd.kr2f_f2.lo248}{5.79}
\wmdef{bd.kr2f_f2.hi248}{5.96}
\wmdef{bd.kr2f_f.def}{0}
\wmdef{bd.kr2f_f.lo}{0}
\wmdef{bd.kr2f_f.hi}{1.8{\times}10^{-10}}
\wmdef{bd.kr2f_f.lo248}{5.79}
\wmdef{bd.kr2f_f.hi248}{5.93}
\wmdef{bd.ii_yield.def}{1}
\wmdef{bd.ii_yield.lo}{0.5}
\wmdef{bd.ii_yield.hi}{1}
\wmdef{bd.ii_yield.lo248}{5.68}
\wmdef{bd.ii_yield.hi248}{5.79}
\wmdef{bd.ii_f_c.def}{0}
\wmdef{bd.ii_f_c.lo}{0}
\wmdef{bd.ii_f_c.hi}{0.5}
\wmdef{bd.ii_f_c.lo248}{5.79}
\wmdef{bd.ii_f_c.hi248}{5.76}
\wmdef{bd.kps_kr.def}{1.7{\times}10^{-10}}
\wmdef{bd.kps_kr.lo}{3.0{\times}10^{-11}}
\wmdef{bd.kps_kr.hi}{1.7{\times}10^{-10}}
\wmdef{bd.kps_kr.lo248}{5.79}
\wmdef{bd.kps_kr.hi248}{5.79}
\wmdef{bd.k5p_f2.def}{1.3{\times}10^{-9}}
\wmdef{bd.k5p_f2.lo}{0}
\wmdef{bd.k5p_f2.hi}{1.3{\times}10^{-9}}
\wmdef{bd.k5p_f2.lo248}{5.76}
\wmdef{bd.k5p_f2.hi248}{5.79}
\wmdef{bd.kr2_f2.def}{2.1{\times}10^{-10}}
\wmdef{bd.kr2_f2.lo}{2.1{\times}10^{-10}}
\wmdef{bd.kr2_f2.hi}{5.0{\times}10^{-10}}
\wmdef{bd.kr2_f2.lo248}{5.79}
\wmdef{bd.kr2_f2.hi248}{5.81}
\wmdef{bd.kr2_pen.def}{0}
\wmdef{bd.kr2_pen.lo}{0}
\wmdef{bd.kr2_pen.hi}{5.0{\times}10^{-10}}
\wmdef{bd.kr2_pen.lo248}{5.79}
\wmdef{bd.kr2_pen.hi248}{5.80}
\wmdef{bd.ib_scale.def}{1}
\wmdef{bd.ib_scale.lo}{0.9}
\wmdef{bd.ib_scale.hi}{1.1}
\wmdef{bd.ib_scale.lo248}{5.79}
\wmdef{bd.ib_scale.hi248}{5.79}
\wmdef{bd.minor}{0.66}
\wmdef{bd.n}{17}
\wmdef{bd.nzero}{7}
\def\wmBands{the trimer's $(T_g/300)^n$ & $-3.32$ & $-3.32$ to $-0.50$ & 5.79 to 2.29 & 4.5 to 4.3 & 4.12 to 1.26 & 0.63 to 0.96 & 742 to 754\\
KrF(B,C) + F \ra Kr + 2F & $0$ & $0$ to $7.8{\times}10^{-10}$\unit{cm^3\,s^{-1}} & 5.79 to 3.99 & 4.5 to 4.4 & 4.12 to 2.42 & 0.63 to 0.48 & 742 to 761\\
\Krts{} + F \ra KrF(B) + Kr & $0$ & $0$ to $1.0{\times}10^{-9}$\unit{cm^3\,s^{-1}} & 5.79 to 7.59 & 4.5 to 0.2 & 4.12 to 5.88 & 0.63 to 0.64 & 742 to 757\\
\el{} + \Krs{} \ra \el{} + \Krpp{}, pre-exponential & $9.0{\times}10^{-7}$ & $0$ to $9.0{\times}10^{-7}$\unit{cm^3\,s^{-1}} & 4.67 to 5.79 & 20.0 to 4.5 & 2.56 to 4.12 & 0.80 to 0.63 & 649 to 742\\
\el{} + \Krs{} \ra \Krp{} + 2\el{}, pre-exponential & $0$ & $0$ to $1.5{\times}10^{-7}$\unit{cm^3\,s^{-1}} & 5.79 to 5.13 & 4.5 to 2.3 & 4.12 to 3.71 & 0.63 to 0.55 & 742 to 767\\
\Krs{} + 2Kr \ra \Krts{} + Kr, 300\unit{K} & $1.7{\times}10^{-32}$ & $1.7{\times}10^{-32}$ to $5.5{\times}10^{-32}$\unit{cm^6\,s^{-1}} & 5.79 to 5.31 & 4.5 to 10.6 & 4.12 to 3.80 & 0.63 to 0.68 & 742 to 706\\
\el{} + \Krpp{} \ra \Krp{} + 2\el{}, pre-exponential & $0$ & $0$ to $2.0{\times}10^{-7}$\unit{cm^3\,s^{-1}} & 5.79 to 5.56 & 4.5 to 3.9 & 4.12 to 3.96 & 0.63 to 0.60 & 742 to 749\\
\Krs{} + \Krs{} \ra \Krp{} + Kr + \el{} & $5.0{\times}10^{-10}$ & $1.5{\times}10^{-10}$ to $5.0{\times}10^{-10}$\unit{cm^3\,s^{-1}} & 6.01 to 5.79 & 5.6 to 4.5 & 4.34 to 4.12 & 0.65 to 0.63 & 732 to 742\\
\KtF{} + \Ft{} & $1.2{\times}10^{-10}$ & $1.2{\times}10^{-10}$ to $4.3{\times}10^{-10}$\unit{cm^3\,s^{-1}} & 5.79 to 5.96 & 4.5 to 4.5 & 4.12 to 4.22 & 0.63 to 0.39 & 742 to 754\\
\KtF{} + F \ra 2Kr + 2F & $0$ & $0$ to $1.8{\times}10^{-10}$\unit{cm^3\,s^{-1}} & 5.79 to 5.93 & 4.5 to 4.5 & 4.12 to 4.24 & 0.63 to 0.26 & 742 to 757\\
KrF\ensuremath{^*} per ion-ion recombination & $1$ & $0.5$ to $1$ & 5.68 to 5.79 & 4.5 to 4.5 & 4.04 to 4.12 & 0.62 to 0.63 & 743 to 742\\
share of the ion-ion KrF\ensuremath{^*} in C & $0$ & $0$ to $0.5$ & 5.79 to 5.76 & 4.5 to 4.5 & 4.12 to 4.11 & 0.63 to 0.63 & 742 to 742\\
\Krpp{} + Kr \ra \Krs{} + Kr & $1.7{\times}10^{-10}$ & $3.0{\times}10^{-11}$ to $1.7{\times}10^{-10}$\unit{cm^3\,s^{-1}} & 5.79 to 5.79 & 5.5 to 4.5 & 4.09 to 4.12 & 1.05 to 0.63 & 727 to 742\\
\Krpp{} + \Ft{}, all channels & $1.3{\times}10^{-9}$ & $0$ to $1.3{\times}10^{-9}$\unit{cm^3\,s^{-1}} & 5.76 to 5.79 & 4.4 to 4.5 & 4.11 to 4.12 & 0.62 to 0.63 & 743 to 742\\
\Krts{} + \Ft{}, both channels & $2.1{\times}10^{-10}$ & $2.1{\times}10^{-10}$ to $5.0{\times}10^{-10}$\unit{cm^3\,s^{-1}} & 5.79 to 5.81 & 4.5 to 3.3 & 4.12 to 4.11 & 0.63 to 0.66 & 742 to 747\\
\Krts{} + \Krts{} \ra \Krtp{} + 2Kr + \el{} & $0$ & $0$ to $5.0{\times}10^{-10}$\unit{cm^3\,s^{-1}} & 5.79 to 5.80 & 4.5 to 3.4 & 4.12 to 4.10 & 0.63 to 0.63 & 742 to 749\\
the wave's electron-ion absorption, a factor & $1$ & $0.9$ to $1.1$ & 5.79 to 5.79 & 4.5 to 4.5 & 4.12 to 4.12 & 0.63 to 0.63 & 742 to 742\\}

% LIGHT v0_4 accuracy numbers (fast). GENERATED by make_overview_v0_4.py
% from LIGHT v0_4 (0.1.4) and ML/. Do not edit. Rerun the script instead.
\providecommand\wfdef[2]{\expandafter\def\csname wf@#1\endcsname{#2}}
\wfdef{sc.eden_kr}{2.3}
\wfdef{sc.eden_f2}{3.2}
\wfdef{sc.ksu_fig2}{4.2}
\wfdef{sc.ksu_fig1}{6.1}
\wfdef{sc.xu_figs5_7}{11.7}
\wfdef{sc.hakuta_fig3}{34.1}
\wfdef{sc.hakuta_fig2}{19.3}
\wfdef{sc.ys_ir300}{\ensuremath{-}0.5}
\wfdef{sc.ys_ratio_230_300}{\ensuremath{-}8.4}
\wfdef{sc.bs_tau_b}{\ensuremath{-}22.3}
\wfdef{sc.same}{10}
\wfdef{sc.train.lo}{2.3}
\wfdef{sc.train.hi}{6.1}
\wfdef{xu.50}{0.91}
\wfdef{xu.100}{0.97}
\wfdef{tah.g0.meas}{6.39}
\wfdef{tah.alpha.meas}{1.11}
\wfdef{tah.net.meas}{5.28}
\wfdef{tah.pure.meas}{0.67}
\wfdef{tah.eta.meas}{0.27}
\wfdef{tah.isat.meas}{1.4}
\wfdef{tah.default.g0}{5.86}
\wfdef{tah.default.g0.x}{0.92}
\wfdef{tah.default.alpha}{0.82}
\wfdef{tah.default.alpha.x}{0.74}
\wfdef{tah.default.net}{5.04}
\wfdef{tah.default.net.x}{0.95}
\wfdef{tah.default.te}{1.18}
\wfdef{tah.default.ne}{$4.51{\times}10^{14}$}
\wfdef{tah.by}{\KtF{} 0.60, \Fm{} 0.09, \Ft{} 0.06, \Krpp{} 0.03, \Krtp{} 0.03, \Krs{} 0.02}
\wfdef{tah.k2f.share}{73}
\wfdef{tah.comp.alpha}{3.81}
\wfdef{tah.comp.alpha.x}{3.43}
\wfdef{tah.kannari.g0}{5.87}
\wfdef{tah.kannari.g0.x}{0.92}
\wfdef{tah.kannari.alpha}{0.84}
\wfdef{tah.kannari.alpha.x}{0.76}
\wfdef{tah.kannari.net}{5.03}
\wfdef{tah.kannari.net.x}{0.95}
\wfdef{tah.kannari.te}{1.17}
\wfdef{tah.kannari.ne}{$4.61{\times}10^{14}$}
\wfdef{tah.nocycle.g0}{5.69}
\wfdef{tah.nocycle.g0.x}{0.89}
\wfdef{tah.nocycle.alpha}{0.84}
\wfdef{tah.nocycle.alpha.x}{0.76}
\wfdef{tah.nocycle.net}{4.84}
\wfdef{tah.nocycle.net.x}{0.92}
\wfdef{tah.nocycle.te}{1.60}
\wfdef{tah.nocycle.ne}{$6.32{\times}10^{14}$}
\wfdef{tah.kpslo.g0}{5.83}
\wfdef{tah.kpslo.g0.x}{0.91}
\wfdef{tah.kpslo.alpha}{0.93}
\wfdef{tah.kpslo.alpha.x}{0.84}
\wfdef{tah.kpslo.net}{4.89}
\wfdef{tah.kpslo.net.x}{0.93}
\wfdef{tah.kpslo.te}{1.26}
\wfdef{tah.kpslo.ne}{$4.82{\times}10^{14}$}
\wfdef{tah.krff.g0}{4.93}
\wfdef{tah.krff.g0.x}{0.77}
\wfdef{tah.krff.alpha}{0.75}
\wfdef{tah.krff.alpha.x}{0.68}
\wfdef{tah.krff.net}{4.18}
\wfdef{tah.krff.net.x}{0.79}
\wfdef{tah.krff.te}{1.18}
\wfdef{tah.krff.ne}{$4.5{\times}10^{14}$}
\wfdef{tah.kr2f.g0}{6.21}
\wfdef{tah.kr2f.g0.x}{0.97}
\wfdef{tah.kr2f.alpha}{0.82}
\wfdef{tah.kr2f.alpha.x}{0.74}
\wfdef{tah.kr2f.net}{5.39}
\wfdef{tah.kr2f.net.x}{1.02}
\wfdef{tah.kr2f.te}{1.19}
\wfdef{tah.kr2f.ne}{$4.49{\times}10^{14}$}
\wfdef{tah.tri.g0}{2.81}
\wfdef{tah.tri.g0.x}{0.44}
\wfdef{tah.tri.alpha}{1.03}
\wfdef{tah.tri.alpha.x}{0.93}
\wfdef{tah.tri.net}{1.78}
\wfdef{tah.tri.net.x}{0.34}
\wfdef{tah.tri.te}{1.19}
\wfdef{tah.tri.ne}{$4.48{\times}10^{14}$}
\wfdef{tah.iiyield.g0}{4.40}
\wfdef{tah.iiyield.g0.x}{0.69}
\wfdef{tah.iiyield.alpha}{0.66}
\wfdef{tah.iiyield.alpha.x}{0.59}
\wfdef{tah.iiyield.net}{3.74}
\wfdef{tah.iiyield.net.x}{0.71}
\wfdef{tah.iiyield.te}{1.19}
\wfdef{tah.iiyield.ne}{$4.49{\times}10^{14}$}
\wfdef{tah.iifc.g0}{5.34}
\wfdef{tah.iifc.g0.x}{0.84}
\wfdef{tah.iifc.alpha}{0.80}
\wfdef{tah.iifc.alpha.x}{0.72}
\wfdef{tah.iifc.net}{4.53}
\wfdef{tah.iifc.net.x}{0.86}
\wfdef{tah.iifc.te}{1.18}
\wfdef{tah.iifc.ne}{$4.52{\times}10^{14}$}
\wfdef{tah.pure.alpha}{0.24}
\wfdef{tah.pure.x}{0.36}
\wfdef{tah.pure.by}{\Krs{} 0.11, \Krpp{} 0.10, \Krtp{} 0.04}
\wfdef{tah.v03.g0}{5.68}
\wfdef{tah.v03.g0.x}{0.89}
\wfdef{tah.v03.alpha}{0.84}
\wfdef{tah.v03.alpha.x}{0.75}
\wfdef{tah.v03.net}{4.85}
\wfdef{tah.v03.net.x}{0.92}
\wfdef{tah.v03.te}{1.60}
\wfdef{tah.v03.ne}{$6.33{\times}10^{14}$}
\wfdef{tah.isat}{1.47}
\wfdef{tah.isat.x}{1.05}
\wfdef{tah.eta}{0.260}
\wfdef{tah.eta.x}{0.96}
\wfdef{tah.form}{0.286}
\wfdef{tah.ml.g0}{5.86}
\wfdef{wv.E}{1500}
\wfdef{wv.f}{95}
\wfdef{wv.em}{0.024}
\wfdef{wv.given}{0.3}
\wfdef{wv.share}{7.5}
\wfdef{wv.248}{5.50}
\wfdef{wv.closure}{$4.0{\times}10^{-14}$}
\wfdef{wv.band}{0.05}
\wfdef{wv.mid.te}{5.46}
\wfdef{wv.mid.ne}{$1.13{\times}10^{14}$}
\wfdef{wv.mid.ei}{0.03}
\wfdef{wv.mid.ei.share}{0.06}
\wfdef{wv.mid.fL}{0.9816}
\wfdef{wv.mid.nuw}{9.6}
\wfdef{wv.ei.15.ghz}{1.46}
\wfdef{wv.ei.15.thz}{0.33}
\wfdef{wv.ei.16.ghz}{7.74}
\wfdef{wv.ei.16.thz}{2.57}
\wfdef{wv.thz.E}{3}
\wfdef{wv.thz.em}{$1.7{\times}10^{-5}$}
\wfdef{wv.err}{10}
\wfdef{wl.gap}{1}
\wfdef{wl.share}{$9.4{\times}10^{-6}$}
\wfdef{wl.ionheat}{$2.0{\times}10^{-6}$}
\wfdef{wl.248}{5.79}
\wfdef{wl.closure}{$2.0{\times}10^{-14}$}
\wfdef{wl.thin.gap}{0.1}
\wfdef{wl.thin.p}{0.1}
\wfdef{wl.thin.share}{0.069}
\wfdef{wl.K0.Krp}{0.84}
\wfdef{wl.K0.Kr2p}{1.22}
\wfdef{wl.K0.Fm}{2.29}
\wfdef{wl.eTe}{2}
\wfdef{wl.iTe}{1}
\wfdef{wl.nu}{63}
\wfdef{wl.1d}{0.036}
\wfdef{wl.1d.cells}{256}
\wfdef{pop.k}{$2.2{\times}10^{-4}$}
\wfdef{pop.db}{$1.2{\times}10^{-5}$}
\wfdef{pop.cells}{2400}
\wfdef{pop.emax}{40}
\wfdef{pop.sigma0}{1.34{\times}10^{-14}}
\wfdef{hf.atoms}{$6.7{\times}10^{-3}$}
\wfdef{hf.ions}{$1.0{\times}10^{-3}$}
\wfdef{hf.share}{1.49}
\wfdef{hf.nuw}{4.4}
\wfdef{hf.te}{1.5}
\wfdef{hf.ne}{$10^{15}$}
\wfdef{pop.short.t}{60}
\wfdef{pop.short.ratio.lo}{1.003}
\wfdef{pop.short.ratio.hi}{1.007}
\wfdef{pop.short.gap}{$2.4{\times}10^{-7}$}
\wfdef{ii.k3}{4.82{\times}10^{-6}}
\wfdef{ii.kL}{5.66{\times}10^{-6}}
\wfdef{ii.k1}{2.61{\times}10^{-6}}
\wfdef{ii.k1.v03}{3.04{\times}10^{-6}}
\wfdef{ii.kL.v03}{8.2{\times}10^{-6}}
\wfdef{ii.k1b}{2.74{\times}10^{-6}}
\wfdef{ii.kLb}{6.35{\times}10^{-6}}
\wfdef{ib.checks}{219}
\wfdef{tau5p}{29.0}
\wfdef{k.KPS}{1.7{\times}10^{-10}}
\wfdef{k.KPS.lo}{3.0{\times}10^{-11}}
\wfdef{k.K27}{1.3{\times}10^{-9}}
\wfdef{k.76.A}{9.0{\times}10^{-7}}
\wfdef{k.76.B}{0.50}
\wfdef{k.76.C}{1.51}
\wfdef{k.stepS.hi}{1.5{\times}10^{-7}}
\wfdef{k.stepS.kannari}{7.0{\times}10^{-8}}
\wfdef{k.stepP.hi}{2.0{\times}10^{-7}}
\wfdef{k.krff.hi}{7.8{\times}10^{-10}}
\wfdef{sig.B}{2.6{\times}10^{-16}}
\wfdef{sig.K2F}{1.0{\times}10^{-18}}
\wfdef{sig.K2F.comp}{6.0{\times}10^{-18}}
\wfdef{sig.Fm}{5.6{\times}10^{-18}}
\wfdef{sig.Kr2p}{1.6{\times}10^{-18}}
\wfdef{sig.Kpp}{6.0{\times}10^{-18}}
\wfdef{sig.F2}{1.2{\times}10^{-20}}
\wfdef{sig.Krs}{1.0{\times}10^{-19}}
\wfdef{n.abs}{6}
\wfdef{n.bands}{17}
\wfdef{n.zero}{7}
\wfdef{ii.yield.lo}{0.5}
\wfdef{ii.fc.hi}{0.5}
\wfdef{ib.lo}{0.9}
\wfdef{ib.hi}{1.1}
\wfdef{K0.Krp}{0.84}
\wfdef{K0.Kr2p}{1.22}
\wfdef{K0.Fm}{2.29}
\def\wfTahVariants{every band at its default & \wf{tah.default.g0} (\x{\wf{tah.default.g0.x}}) & \wf{tah.default.alpha} (\x{\wf{tah.default.alpha.x}}) & \wf{tah.default.net} (\x{\wf{tah.default.net.x}}) & \wf{tah.default.te} & \wf{tah.default.ne}\\
the rung's bands at Kannari's values (\texttt{--ladder kannari}) & \wf{tah.kannari.g0} (\x{\wf{tah.kannari.g0.x}}) & \wf{tah.kannari.alpha} (\x{\wf{tah.kannari.alpha.x}}) & \wf{tah.kannari.net} (\x{\wf{tah.kannari.net.x}}) & \wf{tah.kannari.te} & \wf{tah.kannari.ne}\\
the 5s $\leftrightarrow$ 5p cycle off (\texttt{k\_5s5p} = 0) & \wf{tah.nocycle.g0} (\x{\wf{tah.nocycle.g0.x}}) & \wf{tah.nocycle.alpha} (\x{\wf{tah.nocycle.alpha.x}}) & \wf{tah.nocycle.net} (\x{\wf{tah.nocycle.net.x}}) & \wf{tah.nocycle.te} & \wf{tah.nocycle.ne}\\
\Krpp{} + Kr at the band's low end & \wf{tah.kpslo.g0} (\x{\wf{tah.kpslo.g0.x}}) & \wf{tah.kpslo.alpha} (\x{\wf{tah.kpslo.alpha.x}}) & \wf{tah.kpslo.net} (\x{\wf{tah.kpslo.net.x}}) & \wf{tah.kpslo.te} & \wf{tah.kpslo.ne}\\
KrF(B,C) + F at Kannari's analogue & \wf{tah.krff.g0} (\x{\wf{tah.krff.g0.x}}) & \wf{tah.krff.alpha} (\x{\wf{tah.krff.alpha.x}}) & \wf{tah.krff.net} (\x{\wf{tah.krff.net.x}}) & \wf{tah.krff.te} & \wf{tah.krff.ne}\\
\Krts{} + F at the top of its band & \wf{tah.kr2f.g0} (\x{\wf{tah.kr2f.g0.x}}) & \wf{tah.kr2f.alpha} (\x{\wf{tah.kr2f.alpha.x}}) & \wf{tah.kr2f.net} (\x{\wf{tah.kr2f.net.x}}) & \wf{tah.kr2f.te} & \wf{tah.kr2f.ne}\\
the trimer on $T_g^{-1/2}$ & \wf{tah.tri.g0} (\x{\wf{tah.tri.g0.x}}) & \wf{tah.tri.alpha} (\x{\wf{tah.tri.alpha.x}}) & \wf{tah.tri.net} (\x{\wf{tah.tri.net.x}}) & \wf{tah.tri.te} & \wf{tah.tri.ne}\\
ion-ion KrF\ensuremath{^*} yield at \wf{ii.yield.lo} & \wf{tah.iiyield.g0} (\x{\wf{tah.iiyield.g0.x}}) & \wf{tah.iiyield.alpha} (\x{\wf{tah.iiyield.alpha.x}}) & \wf{tah.iiyield.net} (\x{\wf{tah.iiyield.net.x}}) & \wf{tah.iiyield.te} & \wf{tah.iiyield.ne}\\
ion-ion C share at \wf{ii.fc.hi} & \wf{tah.iifc.g0} (\x{\wf{tah.iifc.g0.x}}) & \wf{tah.iifc.alpha} (\x{\wf{tah.iifc.alpha.x}}) & \wf{tah.iifc.net} (\x{\wf{tah.iifc.net.x}}) & \wf{tah.iifc.te} & \wf{tah.iifc.ne}\\}

% LIGHT v0_4 accuracy numbers (slow). GENERATED by make_overview_v0_4.py
% from LIGHT v0_4 (0.1.4) and ML/. Do not edit. Rerun the script instead.
\providecommand\wfdef[2]{\expandafter\def\csname wf@#1\endcsname{#2}}
\wfdef{ref.248}{5.68}
\wfdef{ref.272}{0.181}
\wfdef{ref.275}{0.825}
\wfdef{ref.420}{4.28}
\wfdef{ref.147}{5.3}
\wfdef{ref.IR}{0.27}
\wfdef{ref.all}{16.5}
\wfdef{ref.tg}{735}
\wfdef{ref.krstar}{65.4}
\wfdef{ref.krplus}{3.9}
\wfdef{ref.krfb}{24.3}
\wfdef{ref.elastic}{2.0}
\wfdef{ref.heat}{71.9}
\wfdef{ref.te}{2.97}
\wfdef{ref.closure}{$10^{-14}$}
\wfdef{ref.solutions}{414}
\wfdef{ref.on}{277}
\wfdef{ref.en.lo}{1.56}
\wfdef{ref.en.hi}{45}
\wfdef{ref.ratio.lo}{1.001}
\wfdef{ref.ratio.hi}{1.004}
\wfdef{ref.gap}{$8.0{\times}10^{-6}$}
\wfdef{ref.brs76}{78.4}
\wfdef{ref.brpKPS}{96.1}
\wfdef{ref.g0}{3.73}
\wfdef{ref.alpha}{0.68}
\wfdef{ref.net}{3.04}
\wfdef{ref.min}{10}
\wfdef{tah.kin.g0}{5.91}
\wfdef{tah.kin.g0.x}{0.93}
\wfdef{tah.kin.alpha}{0.82}
\wfdef{tah.kin.alpha.x}{0.74}
\wfdef{tah.kin.net}{5.09}
\wfdef{tah.kin.net.x}{0.96}
\wfdef{tah.kin.te}{1.16}
\wfdef{tah.kin.ne}{$4.43{\times}10^{14}$}
\wfdef{tah.kin.gap}{$4.6{\times}10^{-4}$}
\wfdef{tah.kin.reached}{every}
\wfdef{tah.kin.en}{0.11}
\wfdef{tah.kin.min}{12}
\wfdef{tah.kin.by}{\KtF{} 0.60, \Fm{} 0.09, \Ft{} 0.06, \Krpp{} 0.03, \Krtp{} 0.03, \Krs{} 0.02}

% LIGHT v0_4 accuracy numbers (uncertainty). GENERATED by make_overview_v0_4.py
% from LIGHT v0_4 (0.1.4) and ML/. Do not edit. Rerun the script instead.
\providecommand\wfdef[2]{\expandafter\def\csname wf@#1\endcsname{#2}}
\wfdef{unc.n}{22}
\wfdef{unc.nband}{17}
\wfdef{unc.nrow}{3}
\wfdef{unc.nsig}{2}
\wfdef{unc.sig.se}{12}
\wfdef{unc.sig.k2f}{20}
\wfdef{unc.sig.k2f.alpha}{15}
\wfdef{unc.min}{6}
\wfdef{unc.ref.y0}{5.79}
\wfdef{unc.ref.quad}{2.28}
\wfdef{unc.ref.quad.pct}{39}
\wfdef{unc.ref.lin}{5.17}
\wfdef{unc.ref.lin.pct}{89}
\wfdef{unc.ref.lo}{-2.17}
\wfdef{unc.ref.hi}{8.17}
\wfdef{unc.ref.down}{-7.96}
\wfdef{unc.ref.up}{+2.39}
\wfdef{unc.ref.top}{the trimer's $(T_g/300)^n$ ($\pm$1.75), \Krts{} + F \ra KrF(B) + Kr ($\pm$0.90), KrF(B,C) + F \ra Kr + 2F ($\pm$0.90), \el{} + \Krs{} \ra \el{} + \Krpp{}, pre-exponential ($\pm$0.56)}
\wfdef{unc.ref.top1}{the trimer's $(T_g/300)^n$}
\wfdef{unc.ref.top1.pct}{30}
\wfdef{unc.g0.y0}{5.86}
\wfdef{unc.g0.quad}{1.93}
\wfdef{unc.g0.quad.pct}{33}
\wfdef{unc.g0.lin}{4.38}
\wfdef{unc.g0.lin.pct}{75}
\wfdef{unc.g0.lo}{-1.39}
\wfdef{unc.g0.hi}{7.37}
\wfdef{unc.g0.down}{-7.25}
\wfdef{unc.g0.up}{+1.51}
\wfdef{unc.g0.top}{the trimer's $(T_g/300)^n$ ($\pm$1.52), KrF\ensuremath{^*} per ion-ion recombination ($\pm$0.73), $\sigma_\mathrm{se}$, 248-nm stimulated emission ($\pm$0.70), KrF(B,C) + F \ra Kr + 2F ($\pm$0.46)}
\wfdef{unc.g0.top1}{the trimer's $(T_g/300)^n$}
\wfdef{unc.g0.top1.pct}{26}
\wfdef{unc.g0.dev}{-0.53}
\wfdef{unc.g0.dev.pct}{-8}
\wfdef{unc.g0.dev.quad}{0.3}
\wfdef{unc.g0.covers}{within}
\wfdef{unc.alpha.y0}{0.82}
\wfdef{unc.alpha.quad}{0.31}
\wfdef{unc.alpha.quad.pct}{37}
\wfdef{unc.alpha.lin}{0.78}
\wfdef{unc.alpha.lin.pct}{96}
\wfdef{unc.alpha.lo}{-0.24}
\wfdef{unc.alpha.hi}{1.33}
\wfdef{unc.alpha.down}{-1.06}
\wfdef{unc.alpha.up}{+0.51}
\wfdef{unc.alpha.top}{\KtF{} + \Ft{} ($\pm$0.18), \KtF{} + F \ra 2Kr + 2F ($\pm$0.16), $\sigma$(\KtF{}) at 248\unit{nm} ($\pm$0.12), the trimer's $(T_g/300)^n$ ($\pm$0.10)}
\wfdef{unc.alpha.top1}{\KtF{} + \Ft{}}
\wfdef{unc.alpha.top1.pct}{22}
\wfdef{unc.alpha.dev}{-0.29}
\wfdef{unc.alpha.dev.pct}{-26}
\wfdef{unc.alpha.dev.quad}{1.0}
\wfdef{unc.alpha.covers}{within}
\wfdef{unc.net.y0}{5.04}
\wfdef{unc.net.quad}{2.01}
\wfdef{unc.net.quad.pct}{40}
\wfdef{unc.net.lin}{4.90}
\wfdef{unc.net.lin.pct}{97}
\wfdef{unc.net.lo}{-2.45}
\wfdef{unc.net.hi}{7.35}
\wfdef{unc.net.down}{-7.49}
\wfdef{unc.net.up}{+2.30}
\wfdef{unc.net.top}{the trimer's $(T_g/300)^n$ ($\pm$1.63), $\sigma_\mathrm{se}$, 248-nm stimulated emission ($\pm$0.70), KrF\ensuremath{^*} per ion-ion recombination ($\pm$0.65), KrF(B,C) + F \ra Kr + 2F ($\pm$0.43)}
\wfdef{unc.net.top1}{the trimer's $(T_g/300)^n$}
\wfdef{unc.net.top1.pct}{32}
\wfdef{unc.net.dev}{-0.24}
\wfdef{unc.net.dev.pct}{-5}
\wfdef{unc.net.dev.quad}{0.1}
\wfdef{unc.net.covers}{within}
\def\wfUncTable{\Krts{} + F \ra KrF(B) + Kr & $0$ to $1.0{\times}10^{-9}$\unit{cm^3\,s^{-1}} & 5.79 to 7.59 & 5.86 to 6.21 & 0.82 to 0.82 & 5.04 to 5.39\\
\Krts{} + \Krts{} \ra \Krtp{} + 2Kr + \el{} & $0$ to $5.0{\times}10^{-10}$\unit{cm^3\,s^{-1}} & 5.79 to 5.80 & 5.86 to 5.86 & 0.82 to 0.82 & 5.04 to 5.04\\
\KtF{} + F \ra 2Kr + 2F & $0$ to $1.8{\times}10^{-10}$\unit{cm^3\,s^{-1}} & 5.79 to 5.93 & 5.86 to 6.00 & 0.82 to 0.50 & 5.04 to 5.50\\
\Krs{} + 2Kr \ra \Krts{} + Kr, 300\unit{K} & $1.7{\times}10^{-32}$ to $5.5{\times}10^{-32}$\unit{cm^6\,s^{-1}} & 5.79 to 5.31 & 5.86 to 5.50 & 0.82 to 0.82 & 5.04 to 4.68\\
\Krts{} + \Ft{}, both channels & $2.1{\times}10^{-10}$ to $5.0{\times}10^{-10}$\unit{cm^3\,s^{-1}} & 5.79 to 5.81 & 5.86 to 5.91 & 0.82 to 0.83 & 5.04 to 5.08\\
\Krs{} + \Krs{} \ra \Krp{} + Kr + \el{} & $1.5{\times}10^{-10}$ to $5.0{\times}10^{-10}$\unit{cm^3\,s^{-1}} & 6.01 to 5.79 & 5.88 to 5.86 & 0.82 to 0.82 & 5.06 to 5.04\\
\KtF{} + \Ft{} & $1.2{\times}10^{-10}$ to $4.3{\times}10^{-10}$\unit{cm^3\,s^{-1}} & 5.79 to 5.96 & 5.86 to 6.05 & 0.82 to 0.46 & 5.04 to 5.59\\
the trimer's $(T_g/300)^n$ & -3.32 to -0.5 & 5.79 to 2.29 & 5.86 to 2.81 & 0.82 to 1.03 & 5.04 to 1.78\\
KrF(B,C) + F \ra Kr + 2F & $0$ to $7.8{\times}10^{-10}$\unit{cm^3\,s^{-1}} & 5.79 to 3.99 & 5.86 to 4.93 & 0.82 to 0.75 & 5.04 to 4.18\\
\Krpp{} + Kr \ra \Krs{} + Kr & $3.0{\times}10^{-11}$ to $1.7{\times}10^{-10}$\unit{cm^3\,s^{-1}} & 5.79 to 5.79 & 5.83 to 5.86 & 0.93 to 0.82 & 4.89 to 5.04\\
\el{} + \Krs{} \ra \el{} + \Krpp{}, pre-exponential & $0$ to $9.0{\times}10^{-7}$\unit{cm^3\,s^{-1}} & 4.67 to 5.79 & 5.69 to 5.86 & 0.84 to 0.82 & 4.84 to 5.04\\
\el{} + \Krs{} \ra \Krp{} + 2\el{}, pre-exponential & $0$ to $1.5{\times}10^{-7}$\unit{cm^3\,s^{-1}} & 5.79 to 5.13 & 5.86 to 5.89 & 0.82 to 0.82 & 5.04 to 5.07\\
\el{} + \Krpp{} \ra \Krp{} + 2\el{}, pre-exponential & $0$ to $2.0{\times}10^{-7}$\unit{cm^3\,s^{-1}} & 5.79 to 5.56 & 5.86 to 5.87 & 0.82 to 0.82 & 5.04 to 5.05\\
\Krpp{} + \Ft{}, all channels & $0$ to $1.3{\times}10^{-9}$\unit{cm^3\,s^{-1}} & 5.76 to 5.79 & 5.85 to 5.86 & 0.82 to 0.82 & 5.03 to 5.04\\
KrF\ensuremath{^*} per ion-ion recombination & 0.5 to 1 & 5.68 to 5.79 & 4.40 to 5.86 & 0.66 to 0.82 & 3.74 to 5.04\\
share of the ion-ion KrF\ensuremath{^*} in C & 0 to 0.5 & 5.79 to 5.76 & 5.86 to 5.34 & 0.82 to 0.80 & 5.04 to 4.53\\
the wave's electron-ion absorption, a factor & 0.9 to 1.1 & 5.79 to 5.79 & 5.86 to 5.86 & 0.82 to 0.82 & 5.04 to 5.04\\
\Krtp{} + \el{} \ra \Krpp{} + Kr & $\pm10$\,\% (Shiu and Biondi 1977, $\pm10$\,\%) & 5.78 to 5.79 & 5.87 to 5.85 & 0.82 to 0.82 & 5.05 to 5.04\\
\Krp{} + 2Kr \ra \Krtp{} + Kr & $\pm9$\,\% (Gr\"ossl et al.\ 1981, $(2.2\pm0.2)\times10^{-31}$) & 5.78 to 5.79 & 5.87 to 5.85 & 0.82 to 0.82 & 5.05 to 5.03\\
\KtF{} radiative lifetime & $\pm9$\,\% (Xu et al.\ 1993, $220\pm20$\unit{ns}) & 5.79 to 5.78 & 5.86 to 5.86 & 0.84 to 0.80 & 5.02 to 5.06\\
$\sigma_\mathrm{se}$, 248-nm stimulated emission & $\pm12$\,\% (Rokni et al.\ 1978, $(2.5\pm0.3)\times10^{-16}$) & -- & 5.16 to 6.56 & -- & 4.34 to 5.74\\
$\sigma$(\KtF{}) at 248\unit{nm} & $\pm20$\,\% (Hakuta et al.\ 1987, $(1.0\pm0.2)\times10^{-18}$) & -- & -- & 0.70 to 0.94 & 4.92 to 5.16\\}

% LIGHT v0_4 energy ladder. GENERATED by make_energy_ladder.py
% from LIGHT v0_4 (0.1.4). Do not edit. Rerun the script instead.
\providecommand\laddef[2]{\expandafter\def\csname lad@#1\endcsname{#2}}
\laddef{version}{0.1.4}
\laddef{E.B}{5.856}
\laddef{DeB}{5.54}
\laddef{BoverKrF}{5.06}
\laddef{wl.G4}{248}
\laddef{hv.G4}{5.00}
\laddef{heat.G4}{0.06}
\laddef{heat.G3}{3.26}
\laddef{heat.K26}{5.54}
\laddef{heat.G5}{0.58}
\laddef{wl.C3}{420}
\laddef{hv.C3}{2.95}
\laddef{heat.C3}{1.53}
\laddef{gapBC}{120}
\laddef{notdrawn}{G6, VIB1, VIB2, VIB3 and VIB4}
\laddef{toe}{G2, PEN\_KR2, K76r and SUPER}
\laddef{zero}{Q2FF, K26C, K26X, K108C, K108X, K2F\_F, PEN\_KR2, KSTEP\_S, KSTEP\_P, QF and QFC}
\laddef{groups}{\texttt{electron}}
\laddef{pale}{}
\laddef{wl.V}{147}
\laddef{hv.V}{8.43}
\laddef{heat.V}{0.77}
\laddef{E.PP}{11.3}
\laddef{gap5p}{1.39}
\laddef{hv.PRAD}{1.39}
\laddef{heat.KPS}{1.39}
\laddef{heat.C1b}{0.70}
\laddef{heat.R11p}{1.55}
\def\ladpicture{%
\definecolor{ladE}{HTML}{2A78D6}% electron gas
\definecolor{ladX}{HTML}{EDA100}% Kr*
\definecolor{ladB}{HTML}{1BAF7A}% KrF(B), 248 and 272 nm
\definecolor{ladC}{HTML}{008300}% KrF(C), 275 nm
\definecolor{ladT}{HTML}{4A3AA7}% Kr2F*, 420 nm
\definecolor{ladV}{HTML}{A63A83}% Kr2*, 147 nm
\definecolor{ladH}{HTML}{E34948}% heat
\definecolor{ladI}{HTML}{52514E}% ions
\definecolor{ladF}{HTML}{A8A69E}% F atoms
\definecolor{ladInk}{HTML}{0B0B0B}%
\definecolor{ladInk2}{HTML}{52514E}%
\definecolor{ladMuted}{HTML}{898781}%
\tikzset{
  lad grid/.style={draw=ladMuted!22, line width=0.3pt},
  lad axisline/.style={draw=ladMuted, line width=0.5pt},
  lad lev/.style={line width=1.4pt},
  lad pair/.style={line width=0.8pt, dash pattern=on 3pt off 2pt},
  lad tick/.style={line width=0.5pt},
  lad arr/.style={line width=0.7pt, -{Triangle[width=3.2pt,length=3.4pt]}},
  lad hv/.style={line width=1.3pt, -{Triangle[width=4.2pt,length=4.4pt]}},
  lad on/.style={},
  lad opt/.style={dash pattern=on 2.4pt off 1.6pt},
  lad zero/.style={dash pattern=on 0.6pt off 1.4pt, line cap=round},
  lad heat/.style={draw=ladH, line width=0.6pt, decorate,
    decoration={snake, amplitude=0.5mm, segment length=1.4mm,
    post length=0.5mm}},
  lad lab/.style={font=\scriptsize, text=ladInk, inner sep=1pt},
  lad lab2/.style={font=\scriptsize, text=ladInk, inner sep=1pt,
    align=right},
  lad lab2l/.style={font=\scriptsize, text=ladInk, inner sep=1pt,
    align=left},
  lad alab/.style={font=\scriptsize, text=ladInk, inner sep=0.8pt,
    fill=white, fill opacity=0.85, text opacity=1, align=center},
  lad slab/.style={font=\tiny, text=ladInk, inner sep=0.6pt,
    fill=white, fill opacity=0.85, text opacity=1, align=right},
  lad hlab/.style={font=\tiny, text=ladInk2, inner sep=0.8pt,
    fill=white, fill opacity=0.85, text opacity=1},
  lad tlab/.style={font=\tiny, text=ladInk2, inner sep=0.8pt},
  lad axis/.style={font=\scriptsize, text=ladInk2, inner sep=1pt},
  lad head/.style={font=\footnotesize\bfseries, text=ladInk2},
}
\begin{tikzpicture}
\draw[lad grid] (0.950,1.000) -- (16.200,1.000);
\draw[lad grid] (0.950,3.000) -- (16.200,3.000);
\draw[lad grid] (0.950,5.000) -- (16.200,5.000);
\draw[lad grid] (0.950,7.000) -- (16.200,7.000);
\draw[lad grid] (0.950,9.000) -- (16.200,9.000);
\draw[lad grid] (0.950,11.000) -- (16.200,11.000);
\draw[lad grid] (0.950,13.000) -- (16.200,13.000);
\draw[lad grid] (0.950,15.000) -- (16.200,15.000);
\draw[lad grid] (0.950,17.000) -- (16.200,17.000);
\draw[lad axisline] (0.950,0.000) -- (0.950,17.450);
\draw[lad lev,draw=ladInk2] (2.650,3.000) -- (6.300,3.000);
\draw[lad lev,draw=ladX] (2.650,12.915) -- (6.300,12.915);
\draw[lad lev,draw=ladI] (2.650,16.996) -- (6.300,16.996);
\draw[lad lev,draw=ladX] (2.650,14.300) -- (6.300,14.300);
\draw[lad lev,draw=ladV] (5.000,12.200) -- (6.300,12.200);
\draw[lad lev,draw=ladI] (5.000,15.846) -- (6.300,15.846);
\draw[lad tick,draw=ladMuted] (3.730,14.998) -- (4.400,14.998);
\draw[lad tick,draw=ladMuted] (3.910,15.750) -- (4.400,15.750);
\draw[lad tick,draw=ladMuted] (4.100,12.915) -- (4.380,12.915);
\draw[lad tick,draw=ladMuted] (5.740,3.770) -- (6.020,3.770);
\draw[lad pair,draw=ladInk2] (10.200,14.396) -- (12.300,14.396);
\draw[lad pair,draw=ladInk2] (10.200,13.246) -- (12.300,13.246);
\draw[lad pair,draw=ladInk2] (7.000,3.800) -- (12.300,3.800);
\draw[lad lev,draw=ladB] (8.400,8.886) -- (11.050,8.886);
\draw[lad lev,draw=ladC] (8.400,8.811) -- (11.050,8.811);
\draw[lad lev,draw=ladT] (7.300,8.276) -- (8.650,8.276);
\draw[lad tick,draw=ladMuted] (9.210,3.856) -- (9.490,3.856);
\draw[lad tick,draw=ladMuted] (9.610,4.296) -- (9.890,4.296);
\draw[lad tick,draw=ladMuted] (10.010,4.331) -- (10.290,4.331);
\draw[lad tick,draw=ladMuted] (7.840,5.326) -- (8.120,5.326);
\draw[lad lev,draw=ladInk2] (13.900,3.000) -- (15.100,3.000);
\draw[lad lev,draw=ladF] (13.900,3.800) -- (15.100,3.800);
\draw[lad pair,draw=ladInk2] (13.900,4.600) -- (15.100,4.600);
\draw[lad pair,draw=ladInk2] (13.900,1.200) -- (15.100,1.200);
\draw[lad lev,draw=ladI] (13.900,0.400) -- (15.100,0.400);
\draw[lad tick,draw=ladMuted] (14.290,6.160) -- (15.100,6.160);
\draw[lad tick,draw=ladMuted] (14.570,7.340) -- (15.100,7.340);
\draw[lad arr,lad on,draw=ladE] (2.800,3.000) -- node[pos=0.52,rotate=90,lad alab] {C2} (2.800,16.996);
\draw[lad arr,lad on,draw=ladE] (3.080,3.000) -- node[pos=0.40,rotate=90,lad alab] {G1} (3.080,12.915);
\draw[lad arr,lad opt,draw=ladE] (3.360,12.915) -- node[pos=0.78,rotate=90,lad alab] {SUPER} (3.360,3.000);
\draw[lad arr,lad on,draw=ladX] (3.080,12.915) -- node[pos=0.36,rotate=90,lad alab] {G2} (3.080,16.996);
\draw[lad arr,lad opt,draw=ladE] (3.360,12.915) -- node[pos=0.62,rotate=90,lad alab] {KSTEP} (3.360,16.996);
\draw[lad arr,lad on,draw=ladE] (3.620,3.000) -- (3.620,14.300);
\draw[lad arr,lad on,draw=ladE] (3.800,3.000) -- node[pos=0.55,rotate=90,lad alab] {C1a--c} (3.800,14.998);
\draw[lad arr,lad on,draw=ladE] (3.980,3.000) -- (3.980,15.750);
\draw[lad heat] (4.400,15.750) -- (4.400,14.300);
\draw[lad hv,lad on,draw=ladX] (4.240,14.300) -- node[pos=0.50,rotate=90,lad alab] {PRAD, near IR} (4.240,12.915);
\draw[lad arr,lad on,draw=ladE] (4.580,12.915) -- node[pos=0.32,rotate=90,lad alab] {K76} (4.580,14.300);
\draw[lad arr,lad on,draw=ladE] (4.820,14.300) -- node[pos=0.30,rotate=90,lad alab] {K76r} (4.820,12.915);
\draw[lad arr,lad on,draw=ladH] (5.060,14.300) -- (5.060,12.915);
\draw[lad arr,lad zero,draw=ladE] (2.940,12.915) -- node[pos=0.85,rotate=90,lad alab] {KSTEP\_S} (2.940,16.996);
\draw[lad arr,lad zero,draw=ladE] (3.500,14.300) -- node[pos=0.85,rotate=90,lad alab] {KSTEP\_P} (3.500,16.996);
\draw[lad arr,lad on,draw=ladH] (2.700,14.300) -- (2.350,13.450);
\draw[lad arr,lad on,draw=ladH] (2.700,12.915) -- (2.350,12.065);
\draw[lad arr,lad on,draw=ladX] (5.100,12.915) -- (5.100,12.200);
\draw[lad arr,lad on,draw=ladI] (5.250,16.996) -- node[pos=0.50,rotate=90,lad alab] {K15} (5.250,15.846);
\draw[lad arr,lad on,draw=ladI] (5.800,15.846) -- node[pos=0.45,rotate=90,lad alab] {R11p} (5.800,14.300);
\draw[lad arr,lad zero,draw=ladV] (6.120,12.200) -- node[pos=0.62,rotate=90,lad alab] {PEN\_KR2} (6.120,15.846);
\draw[lad hv,lad on,draw=ladV] (5.880,12.200) -- node[pos=0.50,rotate=90,lad alab] {KR2\_RAD, 147 nm} (5.880,3.770);
\draw[lad heat] (5.880,3.770) -- (5.880,3.000);
\draw[lad arr,lad on,draw=ladV] (6.300,12.200) -- node[pos=0.42,sloped,lad alab] {K2F\_BR} (7.620,8.276);
\draw[lad arr,lad on,draw=ladV] (6.300,12.200) -- node[pos=0.42,sloped,lad alab] {NK7} (8.660,8.856);
\draw[lad arr,lad zero,draw=ladV] (6.300,12.200) -- node[pos=0.62,sloped,lad alab] {K2F\_F} (8.900,8.856);
\draw[lad arr,lad on,draw=ladX] (6.300,12.915) -- node[pos=0.40,sloped,lad alab] {G3, G3C} (9.140,8.856);
\draw[lad arr,lad on,draw=ladX] (6.300,14.300) -- node[pos=0.30,sloped,lad alab] {K27, K27C} (9.260,8.856);
\draw[lad arr,lad on,draw=ladI] (10.450,13.246) -- node[pos=0.50,rotate=90,lad alab] {K108, K108C} (10.450,8.856);
\draw[lad arr,lad on,draw=ladI] (10.720,14.396) -- node[pos=0.62,rotate=90,lad alab] {K26, K26C} (10.720,8.856);
\draw[lad arr,lad zero,draw=ladH] (10.300,13.246) -- (10.300,12.446);
\draw[lad arr,lad zero,draw=ladH] (10.300,14.396) -- (10.300,13.796);
\draw[lad arr,lad on,draw=ladB] (8.520,8.856) -- (8.520,8.276);
\draw[lad hv,lad on,draw=ladB] (9.350,8.856) -- node[pos=0.50,rotate=90,lad alab] {G4, 248 nm} (9.350,3.856);
\draw[lad hv,lad on,draw=ladB] (9.750,8.856) -- node[pos=0.50,rotate=90,lad alab] {G4A, 272 nm} (9.750,4.296);
\draw[lad heat] (9.750,4.296) -- (9.750,3.800);
\draw[lad hv,lad on,draw=ladC] (10.150,8.841) -- node[pos=0.50,rotate=90,lad alab] {G4C, 275 nm} (10.150,4.331);
\draw[lad heat] (10.150,4.331) -- (10.150,3.800);
\draw[lad arr,lad on,draw=ladH] (10.930,8.841) -- (10.930,8.041);
\draw[lad hv,lad on,draw=ladT] (7.980,8.276) -- node[pos=0.64,rotate=90,lad alab] {C3, 420 nm} (7.980,5.326);
\draw[lad heat] (7.980,5.326) -- (7.980,3.800);
\draw[lad arr,lad on,draw=ladH] (7.520,8.276) -- (7.520,7.476);
\draw[lad arr,lad zero,draw=ladH] (8.400,8.276) -- (8.400,7.476);
\draw[lad arr,lad on,draw=ladE] (14.360,3.000) -- node[pos=0.62,rotate=90,lad alab] {DIS\_a3} (14.360,6.160);
\draw[lad arr,lad on,draw=ladE] (14.640,3.000) -- node[pos=0.66,rotate=90,lad alab] {DIS\_A1} (14.640,7.340);
\draw[lad heat] (15.100,7.340) -- (15.100,4.600);
\draw[lad arr,lad on,draw=ladE] (14.080,3.000) -- node[pos=0.50,rotate=90,lad alab] {ATT} (14.080,1.200);
\draw[lad arr,lad on,draw=ladF] (14.880,4.600) -- node[pos=0.25,rotate=90,lad alab] {FF} (14.880,3.000);
\node[lad axis,anchor=east] at (0.870,1.000) {\ensuremath{-}2};
\node[lad axis,anchor=east] at (0.870,3.000) {0};
\node[lad axis,anchor=east] at (0.870,5.000) {2};
\node[lad axis,anchor=east] at (0.870,7.000) {4};
\node[lad axis,anchor=east] at (0.870,9.000) {6};
\node[lad axis,anchor=east] at (0.870,11.000) {8};
\node[lad axis,anchor=east] at (0.870,13.000) {10};
\node[lad axis,anchor=east] at (0.870,15.000) {12};
\node[lad axis,anchor=east] at (0.870,17.000) {14};
\node[lad axis,anchor=east] at (0.870,18.000) {eV};
\node[lad head,anchor=center] at (4.475,18.000) {Kr and Kr\ensuremath{_2}};
\node[lad head,anchor=center] at (9.650,18.000) {KrF and Kr\ensuremath{_2}F};
\node[lad head,anchor=center] at (14.500,18.000) {F\ensuremath{_2}};
\node[lad lab,anchor=east] at (2.570,3.000) {Kr};
\node[lad lab,anchor=east] at (2.570,12.915) {Kr\ensuremath{^*} 9.915};
\node[lad lab,anchor=east] at (2.570,16.996) {Kr\ensuremath{^+} 13.996};
\node[lad lab,anchor=east] at (2.570,14.300) {Kr\ensuremath{^{**}} 11.3};
\node[lad lab2l,anchor=north west] at (5.030,12.150) {Kr\ensuremath{_2^*}\\9.2};
\node[lad lab2,anchor=south east] at (6.300,15.896) {Kr\ensuremath{_2^+}\\12.846};
\node[lad tlab,anchor=west] at (4.460,14.998) {11.998};
\node[lad tlab,anchor=west] at (4.460,15.750) {12.75};
\node[lad hlab,anchor=west] at (4.340,12.615) {0.00};
\node[lad alab,anchor=west] at (5.120,13.607) {KPS};
\node[lad alab,anchor=east] at (2.310,13.580) {K27X};
\node[lad alab,anchor=east] at (2.310,12.195) {G3X};
\node[lad alab,anchor=east] at (5.040,12.557) {AK8\_Kr};
\node[lad hlab,anchor=west] at (5.980,3.370) {0.77};
\node[lad lab,anchor=west] at (12.380,14.396) {Kr\ensuremath{^+} + F\ensuremath{^-} 11.396};
\node[lad lab,anchor=west] at (12.380,13.246) {Kr\ensuremath{_2^+} + F\ensuremath{^-} 10.246};
\node[lad lab,anchor=south east] at (12.300,3.840) {Kr + F 0.8};
\node[lad lab2l,anchor=west] at (11.130,8.856) {KrF(B) 5.856\\KrF(C) 120\,cm\ensuremath{^{-1}} below\\BC, CB mix them};
\node[lad lab2,anchor=east] at (7.220,8.276) {Kr\ensuremath{_2}F\ensuremath{^*}\\5.276};
\node[lad alab,anchor=east] at (10.240,12.846) {K108X};
\node[lad alab,anchor=east] at (10.240,14.096) {K26X};
\node[lad slab,anchor=east] at (8.470,8.566) {G5,\\G5C};
\node[lad hlab,anchor=north] at (9.350,3.420) {X, 0.06};
\node[lad hlab,anchor=north] at (9.750,3.720) {A, 0.50};
\node[lad hlab,anchor=north] at (10.150,3.420) {A, 0.53};
\node[lad alab,anchor=north west] at (10.620,7.991) {QF2, QKR, QE, QF\\QF2C, QKRC, QEC, QFC};
\node[lad hlab,anchor=north] at (7.980,3.720) {1.53};
\node[lad alab,anchor=north] at (7.520,7.426) {Q2F};
\node[lad alab,anchor=north] at (8.400,7.426) {Q2FF};
\node[lad lab,anchor=east] at (13.820,3.000) {F\ensuremath{_2}};
\node[lad lab,anchor=east] at (13.820,3.800) {F 0.8};
\node[lad lab,anchor=east] at (13.820,4.600) {F + F 1.6};
\node[lad lab,anchor=east] at (13.820,1.200) {F\ensuremath{^-} + F \ensuremath{-}1.8};
\node[lad lab,anchor=east] at (13.820,0.400) {F\ensuremath{^-} \ensuremath{-}2.6};
\node[lad tlab,anchor=west] at (15.150,6.160) {{\unsansmath\ensuremath{a\,^3\Pi_u}} 3.16};
\node[lad tlab,anchor=west] at (15.150,7.340) {{\unsansmath\ensuremath{A\,^1\Pi_u}} 4.34};
\draw[lad arr,lad on,draw=ladInk2] (1.250,-0.550) -- (1.950,-0.550);
\node[lad lab,anchor=west] at (2.050,-0.550) {a row of the default network};
\draw[lad arr,lad opt,draw=ladInk2] (8.850,-0.550) -- (9.550,-0.550);
\node[lad lab,anchor=west] at (9.650,-0.550) {a row of an optional group, off by default};
\draw[lad arr,lad zero,draw=ladInk2] (1.250,-0.970) -- (1.950,-0.970);
\node[lad lab,anchor=west] at (2.050,-0.970) {a band whose default is 0, no measurement};
\draw[lad hv,lad on,draw=ladB] (8.850,-0.970) -- (9.550,-0.970);
\node[lad lab,anchor=west] at (9.650,-0.970) {a photon, with its row and band};
\draw[lad arr,lad on,draw=ladH] (1.250,-1.390) -- (1.950,-1.390);
\node[lad lab,anchor=west] at (2.050,-1.390) {a quench, the level's energy to heat};
\draw[lad heat] (8.850,-1.390) -- (9.550,-1.390);
\node[lad lab,anchor=west] at (9.650,-1.390) {heat as a lower state falls apart (eV) or cascades};
\draw[lad pair,draw=ladInk2] (1.250,-1.810) -- (1.950,-1.810);
\node[lad lab,anchor=west] at (2.050,-1.810) {a pair, the sum of two formation energies};
\draw[lad tick,draw=ladMuted] (8.850,-1.810) -- (9.550,-1.810);
\node[lad lab,anchor=west] at (9.650,-1.810) {a threshold, or where a photon leaves its lower state};
\end{tikzpicture}}

\makeatletter
\newcommand\wi[1]{\ifcsname wi@#1\endcsname\csname wi@#1\endcsname
  \else\errmessage{no v0_2 number '#1'. Rerun make_overview_v0_2.py}\fi}
\newcommand\wk[1]{\ifcsname wk@#1\endcsname\csname wk@#1\endcsname
  \else\errmessage{no v0_3 number '#1'. Rerun make_overview_v0_3.py}\fi}
\newcommand\wm[1]{\ifcsname wm@#1\endcsname\csname wm@#1\endcsname
  \else\errmessage{no map number '#1'. Rerun make_system_map_v0_4.py}\fi}
\newcommand\wf[1]{\ifcsname wf@#1\endcsname\csname wf@#1\endcsname
  \else\errmessage{no v0_4 number '#1'. Rerun make_overview_v0_4.py}\fi}
\newcommand\lad[1]{\ifcsname lad@#1\endcsname\csname lad@#1\endcsname
  \else\errmessage{no ladder number '#1'. Rerun make_energy_ladder.py}\fi}
\makeatother

% chemistry and units
\newcommand{\unit}[1]{\ensuremath{\,\mathrm{#1}}}
\newcommand{\el}{\ensuremath{\mathrm{e}^-}}
\newcommand{\Krs}{Kr\ensuremath{^*}}
\newcommand{\Krpp}{Kr\ensuremath{^{**}}}
\newcommand{\Krp}{Kr\ensuremath{^+}}
\newcommand{\Krtp}{Kr\ensuremath{_2^+}}
\newcommand{\Krts}{Kr\ensuremath{_2^*}}
\newcommand{\KtF}{Kr\ensuremath{_2}F\ensuremath{^*}}
\newcommand{\Ft}{F\ensuremath{_2}}
\newcommand{\Fm}{F\ensuremath{^-}}
\newcommand{\ra}{\ensuremath{\to}\ }
\newcommand{\Edep}{\ensuremath{E_\mathrm{dep}}}
\newcommand{\gIP}{\ensuremath{g_0I_\mathrm{sat}/P_\mathrm{in}}}
\newcommand{\x}[1]{\ensuremath{\times}#1}
\newcommand{\pc}[1]{\wm{#1}\,\%}
\newcommand{\Lam}{\ensuremath{\text{{\unsansmath$\Lambda$}}}}
\renewcommand{\topfraction}{0.95}
\renewcommand{\textfraction}{0.03}
\newcommand{\cm}[1]{\ensuremath{\,\mathrm{cm^{#1}s^{-1}}}}
\newcommand{\pctof}[2][0]{\pgfmathparse{100*#2}%
  \pgfmathprintnumber[fixed,precision=#1]{\pgfmathresult}}
% a model-over-measurement factor as a signed percent error
\newcommand{\errof}[1]{\pgfmathparse{100*(#1-1)}%
  \pgfmathprintnumber[fixed,precision=0,showpos]{\pgfmathresult}}

\begin{document}
\sansmath

{\LARGE\bfseries How accurate LIGHT v0\_4 is}\par
\vspace{4pt}
{\large The basic network of LIGHT v0\_4, version \lad{version}, on 9
October 2026: v0\_3 with the 5p rung of krypton, the 248-nm optics,
KrF(B,C) + F, ion-ion recombination's limbs and bands, the ions' drift to
the walls and the EM wave, with its energy ladder and its checks}\par
\vspace{10pt}

\noindent\fbox{\parbox{\dimexpr\linewidth-2\fboxsep-2\fboxrule\relax}{%
\setlength{\parskip}{5pt}%
\textbf{Bottom line.} LIGHT v0\_4 reproduces the fitted KrF
decays~\cite{eden78,gadomski92} within \wf{sc.train.hi}\,\% rms
(\wf{sc.train.lo} to \wf{sc.train.hi}\,\%) and the decays no fit
used~\cite{xu93} within \wf{sc.xu_figs5_7}\,\% rms, as v0\_2 and v0\_3 do.
At the Tahoma laser~\cite{tahoma89} it gives the small-signal gain to
\errof{\wf{tah.kin.g0.x}}\,\% (\wf{tah.kin.g0} against the measured
\wf{tah.g0.meas}\unit{\%/cm}), the absorption, scored for the first time,
to \errof{\wf{tah.kin.alpha.x}}\,\% (\wf{tah.kin.alpha} against
\wf{tah.alpha.meas}), the net gain to \errof{\wf{tah.kin.net.x}}\,\%, the
saturation intensity to \errof{\wf{tah.isat.x}}\,\% and the formation
efficiency to \errof{\wf{tah.eta.x}}\,\%. No fit used Tahoma's numbers and
no rate was set by a laser measurement.

Its own uncertainty, propagated from its sources (Section~\ref{sec:accuracy}:
\wf{unc.n} inputs, each at the two ends its source allows, combined in
quadrature), is $\pm$\wf{unc.ref.quad.pct}\,\% on the 248-nm share at the
reference point and $\pm$\wf{unc.g0.quad.pct}, $\pm$\wf{unc.alpha.quad.pct}
and $\pm$\wf{unc.net.quad.pct}\,\% on Tahoma's gain, absorption and net
gain. The largest term everywhere but the absorption is
\wf{unc.ref.top1}, a band with no measurement behind it. Every Tahoma
measurement lies within one quadrature unit of the model (the gain
\wf{unc.g0.dev.quad}, the absorption \wf{unc.alpha.dev.quad}, the net gain
\wf{unc.net.dev.quad} units): at Tahoma the model's own uncertainty, not
its error, is the limit.

The largest errors: Hakuta et al.'s C-A decays~\cite{hakuta90},
\wf{sc.hakuta_fig2} and \wf{sc.hakuta_fig3}\,\% rms; Burnham and Searles'
KrF(B) lifetime~\cite{burnham77}, \wf{sc.bs_tau_b}\,\%; the VUV yield of
pure krypton at low pump~\cite{morozov08}, \errof{\wk{pure.vuv.x}}\,\%
(v0\_3's run, not repeated here); and Tahoma's absorption,
\errof{\wf{tah.kin.alpha.x}}\,\%. \KtF{} makes \wf{tah.k2f.share}\,\% of
that absorption at the cross section Hakuta et al.\ measured~\cite{hakuta87};
at the compilation value every KrF model carried, the absorption would be
\x{\wf{tah.comp.alpha.x}} the measurement. What v0\_4 changes at the
reference point is the 5s $\leftrightarrow$ 5p electron cycle on Kannari's
fit~\cite{kannari85,jacob76}: it turns \pc{cyc.e} of the input into heat
and takes the 248-nm share from \wm{om.v03} to \wm{out.248}\,\% of \Edep{}
(\wf{ref.248}\,\% under the kinetic EEDF). That fit has no direct
measurement behind it, and its band runs to 0.}}
\medskip

LIGHT v0\_4 is v0\_3 with what the audit of 8 October 2026 found missing
against the measured data and the modern kinetic codes: the 5p rung of
krypton (\Krpp{}, fed by the higher excitations and by \Krtp{} + \el{},
relaxed by Kr at a measured rate, radiating at the measured manifold
lifetime, harpooning \Ft{} at a measured rate, pumped from \Krs{} and
returned by the electrons on Kannari's fit, its two stepwise ionizations
bands at 0), KrF(B,C) + F as a band, the 248-nm optics from Kannari's
absorber block with Hakuta et al.'s measured \KtF{} cross section, ion-ion
recombination on Flannery's turnover with the Langevin limb of the
measured mobilities and its yield and B/C bands, the ions' drift to the
walls of a gap, and an EM wave absorbed by the Lorentz gas on the atoms
plus Turnbull et al.'s electron-ion inverse bremsstrahlung with their
model's error as a band. The kinetic EEDF carries the 5s $\leftrightarrow$
5p collisions, the ions' Coulomb collisions and the wave's high-frequency
field term. The \wf{sc.same} kinetic datasets have no electrons and no
\Krpp{}, so v0\_4 scores them exactly as v0\_2, and
\texttt{make\_overview\_v0\_4.py} refuses to write a number that differs.
Table~\ref{tab:score} lists every test.

\begin{table}[!tbp]
\centering\footnotesize
\renewcommand{\arraystretch}{1.12}
\begin{tabularx}{\linewidth}{@{}>{\raggedright\arraybackslash}X
  >{\raggedright\arraybackslash}p{2.3cm}>{\raggedright\arraybackslash}p{2.4cm}
  >{\raggedright\arraybackslash}p{2.3cm}>{\raggedright\arraybackslash}p{1.9cm}@{}}
\toprule
test & measured & v0\_4 & v0\_3 & role of the data\\
\midrule
energy ledger closure & & below \wf{wl.closure} with walls,
  \wf{wv.closure} with the wave & below \wk{closure.max} & bookkeeping\\
KrF(B) decay against Kr and F$_2$, Eden 1978 & & rms \wf{sc.eden_kr} and
  \wf{sc.eden_f2}\,\% & the same & train\\
KrF(B,C) slow decay, Gadomski 1992 Figs.\ 1 and 2 & & rms
  \wf{sc.ksu_fig1} and \wf{sc.ksu_fig2}\,\% & the same & train\\
C/B emission ratio at 300\unit{K}, Yu and Setser 1990 & \wi{ys.meas}
  $\pm$ \wi{ys.err} & \wf{sc.ys_ir300}\,\% & the same & train\\
KrF(B,C) decays, Xu 1993 Figs.\ 5 to 7 & & rms \wf{sc.xu_figs5_7}\,\% &
  the same & held out\\
decay constants Xu 1993 printed, 50 and 100 Torr Kr & \wi{xu.50.meas}
  and \wi{xu.100.meas}\unit{s^{-1}} & \x{\wf{xu.50}} and \x{\wf{xu.100}} &
  the same & structure check\\
C-A decays, Hakuta 1990 & & rms \wf{sc.hakuta_fig2} and
  \wf{sc.hakuta_fig3}\,\% & the same & held out, in conflict\\
$\tau_B$, Burnham and Searles 1977 & \wi{sc.bs.meas} $\pm$
  \wi{sc.bs.err}\unit{ns} & \wf{sc.bs_tau_b}\,\% & the same & held out, in
  conflict\\
C/B ratio 230\unit{K} over 300\unit{K}, Yu and Setser & \wi{ysr.meas}
  $\pm$ \wi{ysr.err} & \wf{sc.ys_ratio_230_300}\,\% & the same & held out\\
\midrule
Tahoma small-signal gain $g_0$, kinetic EEDF & \wf{tah.g0.meas}\unit{\%/cm} &
  \wf{tah.kin.g0}, \x{\wf{tah.kin.g0.x}} & \wk{tah.g0}, \x{\wk{tah.g0.x}} &
  device, score only\\
\quad the same, v0\_2's closure & & \wf{tah.default.g0},
  \x{\wf{tah.default.g0.x}} & \wf{tah.v03.g0}, \x{\wf{tah.v03.g0.x}} &
  device, score only\\
Tahoma absorption $\alpha_0$, kinetic EEDF & \wf{tah.alpha.meas}\unit{\%/cm}
  & \wf{tah.kin.alpha}, \x{\wf{tah.kin.alpha.x}} & not computed & device,
  score only\\
\quad the same, v0\_2's closure & & \wf{tah.default.alpha},
  \x{\wf{tah.default.alpha.x}} & \wf{tah.v03.alpha}, \x{\wf{tah.v03.alpha.x}}
  (v0\_4's cross sections on v0\_3's densities) & device, score only\\
Tahoma net gain $g_0-\alpha_0$, kinetic EEDF & \wf{tah.net.meas}\unit{\%/cm}
  & \wf{tah.kin.net}, \x{\wf{tah.kin.net.x}} & & device, score only\\
Tahoma loss in pure Kr, 257.3\unit{nm} & \wf{tah.pure.meas}\unit{\%/cm} &
  \wf{tah.pure.alpha}, \x{\wf{tah.pure.x}} (248-nm cross sections) & &
  rough check\\
Tahoma saturation intensity $I_\mathrm{sat}$, v0\_2's closure &
  \wf{tah.isat.meas} $\pm$ \wi{tah.isat.err}\unit{MW/cm^2} & \wf{tah.isat},
  \x{\wf{tah.isat.x}} & \wk{tah.isat}, \x{\wk{tah.isat.x}} (kinetic) &
  device, score only\\
Tahoma formation efficiency \gIP, v0\_2's closure & \wf{tah.eta.meas} &
  \wf{tah.eta}, \x{\wf{tah.eta.x}} & \wk{tah.eta}, \x{\wk{tah.eta.x}}
  (kinetic) & device, score only\\
\bottomrule
\end{tabularx}
\caption{LIGHT v0\_4 against the data, beside v0\_3, every band at its
default unless a row says otherwise. A percentage is the deviation from
the data, signed for a single value and rms for a set. A \x{} factor is
the model over the measurement. ``The same'' is checked: the generator
stops if v0\_4's number differs from v0\_2's. The absorption rows are new:
Tahoma's Table 3-3 (beam stop in, 330\unit{kW\,cm^{-3}}) prints $g_0$,
$\alpha_0$ and $g_0-\alpha_0$, and Table 3-2 the pure-Kr loss at
257.3\unit{nm}.}
\label{tab:score}
\end{table}

\section{Overall accuracy and its uncertainty}
\label{sec:accuracy}

\textbf{Against the data.} Table~\ref{tab:score} scores every test;
Table~\ref{tab:errors} reads the same numbers as errors, the model's
deviation from the measurement in percent of the measurement (rms over a
set of points), and names the largest.

\begin{table}[H]
\centering\footnotesize
\renewcommand{\arraystretch}{1.1}
\begin{tabularx}{\linewidth}{@{}>{\raggedright\arraybackslash}X
  >{\raggedright\arraybackslash}p{3.0cm}>{\raggedright\arraybackslash}p{3.0cm}
  >{\raggedright\arraybackslash}p{3.4cm}@{}}
\toprule
quantity & measured & v0\_4 & error\\
\midrule
KrF(B,C) decays the fit used, Eden 1978 and Gadomski 1992~\cite{eden78,gadomski92}
  & four sets & & within \wf{sc.train.hi}\,\% rms (\wf{sc.train.lo} to
  \wf{sc.train.hi})\\
KrF(B,C) decays no fit used, Xu 1993~\cite{xu93} & three figures & & within
  \wf{sc.xu_figs5_7}\,\% rms\\
C/B emission ratio at 300\unit{K}~\cite{yu90} & \wi{ys.meas} $\pm$ \wi{ys.err} & &
  \wf{sc.ys_ir300}\,\%\\
its 230\unit{K}/300\unit{K} ratio~\cite{yu90} & \wi{ysr.meas} $\pm$ \wi{ysr.err} & &
  \wf{sc.ys_ratio_230_300}\,\%\\
decay constants, 50 and 100\unit{Torr} Kr~\cite{xu93} & \wi{xu.50.meas},
  \wi{xu.100.meas}\unit{s^{-1}} & \x{\wf{xu.50}}, \x{\wf{xu.100}} &
  \errof{\wf{xu.50}}, \errof{\wf{xu.100}}\,\%\\
Tahoma small-signal gain $g_0$~\cite{tahoma89} & \wf{tah.g0.meas}\unit{\%/cm} &
  \wf{tah.kin.g0} & \errof{\wf{tah.kin.g0.x}}\,\%\\
Tahoma absorption $\alpha_0$ & \wf{tah.alpha.meas}\unit{\%/cm} & \wf{tah.kin.alpha} &
  \errof{\wf{tah.kin.alpha.x}}\,\%\\
Tahoma net gain $g_0-\alpha_0$ & \wf{tah.net.meas}\unit{\%/cm} & \wf{tah.kin.net} &
  \errof{\wf{tah.kin.net.x}}\,\%\\
Tahoma saturation intensity & \wf{tah.isat.meas} $\pm$ \wi{tah.isat.err}\unit{MW/cm^2}
  & \wf{tah.isat} & \errof{\wf{tah.isat.x}}\,\%\\
Tahoma formation efficiency \gIP{} & \wf{tah.eta.meas} & \wf{tah.eta} &
  \errof{\wf{tah.eta.x}}\,\%\\
\midrule
\multicolumn{4}{@{}l}{\emph{the largest errors}}\\
C-A decays, Hakuta 1990~\cite{hakuta90} (in conflict with the fitted sets) & two
  figures & & \wf{sc.hakuta_fig2} and \wf{sc.hakuta_fig3}\,\% rms\\
KrF(B) lifetime, Burnham and Searles 1977~\cite{burnham77} (in conflict) &
  \wi{sc.bs.meas} $\pm$ \wi{sc.bs.err}\unit{ns} & & \wf{sc.bs_tau_b}\,\%\\
VUV yield of pure Kr at low pump, Morozov 2008~\cite{morozov08}, v0\_3's run &
  \wk{pure.vuv.meas} $\pm$ \wk{pure.vuv.err} & \wk{pure.vuv} &
  \errof{\wk{pure.vuv.x}}\,\%\\
Tahoma absorption $\alpha_0$ & \wf{tah.alpha.meas}\unit{\%/cm} & \wf{tah.kin.alpha}
  & \errof{\wf{tah.kin.alpha.x}}\,\%\\
\bottomrule
\end{tabularx}
\caption{The scores of Table~\ref{tab:score} as errors. A single value's
error is signed, model less measurement over the measurement; a set's is
the rms over its points. Tahoma's are under the kinetic EEDF; the rest
have no electrons.}
\label{tab:errors}
\end{table}

\textbf{Propagated from the sources.} The model's own uncertainty is what
its inputs' uncertainties do to its outputs. Each input is set to the two
ends its source allows and the outputs are recomputed (Table~\ref{tab:unc},
under v0\_2's closure; the kinetic one moves these outputs by about 1\,\%):
each of the \wf{unc.nband} bands to its range (a published spread, or 0 to
an estimate where nothing is measured, Table~\ref{tab:bands}); the
\wf{unc.nrow} constants the model uses at a measured value whose source
quotes an error and that no band covers, \Krtp{} + \el{}
($\pm$10\,\%~\cite{shiu77}), \Krp{} + 2Kr ($\pm$9\,\%~\cite{grossl81}) and
the \KtF{} lifetime ($\pm$9\,\%~\cite{xu93}), to $1\pm$ that error; and the
\wf{unc.nsig} cross sections with a quoted error analytically,
$\sigma_\mathrm{se}$ $\pm$\wf{unc.sig.se}\,\%~\cite{tellinghuisen76,kannari85}
on $g_0$ and $\sigma$(\KtF) $\pm$\wf{unc.sig.k2f}\,\%~\cite{hakuta87} on
its share of $\alpha_0$, which is $\pm$\wf{unc.sig.k2f.alpha}\,\% of
$\alpha_0$. Half the difference between an input's two ends is its term.
The quadrature sum treats the terms as independent and each as one
standard deviation, which a band's range is not, so it is an indicative
number. The linear sums of the downward and of the upward excursions,
every input pushing the same way, overestimate the spread (the inputs
draw on the same KrF(B) and a share cannot fall below 0) and are printed
as the ceiling on it, not as a bound. Not propagated: the electron-impact cross sections, which carry no quoted
error, and the fitted KrF(B,C) set (the lifetimes, the B-C mixing, the
quenches by \el, \Ft{} and Kr, the trimer rate at 300\unit{K}), which the
training decays pin at \wf{sc.train.hi}\,\% rms and Table~\ref{tab:errors}
scores.

On the 248-nm share at the reference point (\wf{unc.ref.y0}\,\% of \Edep)
the quadrature sum is $\pm$\wf{unc.ref.quad} points
($\pm$\wf{unc.ref.quad.pct}\,\%), the linear sum $\pm$\wf{unc.ref.lin}
points and the excursions sum to \wf{unc.ref.down} and \wf{unc.ref.up}
points; the largest terms are \wf{unc.ref.top}. On Tahoma's gain
(\wf{unc.g0.y0}\unit{\%/cm}) it is $\pm$\wf{unc.g0.quad.pct}\,\% in
quadrature (\wf{unc.g0.down} and \wf{unc.g0.up}\unit{\%/cm} summed), the
largest terms \wf{unc.g0.top}; the measurement is \wf{unc.g0.dev.pct}\,\%
away, \wf{unc.g0.dev.quad} quadrature units, \wf{unc.g0.covers} one unit.
On the absorption (\wf{unc.alpha.y0}\unit{\%/cm}) it is
$\pm$\wf{unc.alpha.quad.pct}\,\% (\wf{unc.alpha.down} and
\wf{unc.alpha.up} summed), the largest terms \wf{unc.alpha.top}; the
measurement is \wf{unc.alpha.dev.pct}\,\% away, \wf{unc.alpha.dev.quad}
units, \wf{unc.alpha.covers} one unit. On the net gain
(\wf{unc.net.y0}\unit{\%/cm}) it is $\pm$\wf{unc.net.quad.pct}\,\%
(\wf{unc.net.down} and \wf{unc.net.up} summed), the measurement
\wf{unc.net.dev.pct}\,\% away, \wf{unc.net.dev.quad} units. So at Tahoma
the model's own uncertainty is larger than its errors: a third to two
fifths of each output, against errors of a few percent on the gain and
the net gain and \errof{\wf{tah.kin.alpha.x}}\,\% on the absorption, and
one unmeasured law, the trimer's $(T_g/300)^n$, carries most of it on the
gain while the \KtF{} quenches and its cross section carry the
absorption's. Measuring that law, or the trimer rate at the operating
temperature, would tighten the gain's uncertainty more than any other
measurement.

\begin{table}[H]
\centering\scriptsize
\resizebox{\linewidth}{!}{%
\begin{tabular}{@{}llllll@{}}
\toprule
input & its two ends & 248\unit{nm} (\% of \Edep) & $g_0$ (\%/cm) & $\alpha_0$ (\%/cm) & $g_0-\alpha_0$\\
\midrule
\wfUncTable
\bottomrule
\end{tabular}}
\caption{The propagated uncertainty's inputs, each alone at its two ends,
the others at their defaults, and the four outputs at those ends (low end
to high end). The reference point is the notes' (1\unit{atm}, 0.5\,\% \Ft,
300\unit{kW\,cm^{-3}} for 1\unit{\mu s}); Tahoma is the 200 to 400-ns window
under the beam drive. The cross sections enter analytically, so they have
no reference-point column. \wf{unc.min} minutes of runs.}
\label{tab:unc}
\end{table}

\section{What v0\_4 changes}

\textbf{The 5p rung.} One lumped \Krpp{} for the 4p$^5$5p manifold at
\wm{E.PP}\unit{eV}, the Biagi group's own threshold. C1a to C1c make it,
\Krtp{} + \el{} makes it (Shiu and Biondi~\cite{shiu77} saw the 5p lines), Kr relaxes it
to \Krs{} at $\wf{k.KPS}\cm{3}$ (Gazeli et al.~\cite{gazeli21}, the total loss of
one 5p$'$ level in Kr, the band down to the Ar analogue $\wf{k.KPS.lo}$~\cite{jancaitis84}),
it radiates to \Krs{} at the measured manifold lifetime of
\wf{tau5p}\unit{ns} (the statistical mean of the ten measured level
lifetimes~\cite{dzierzega00}), it harpoons \Ft{} at
$\wf{k.K27}\cm{3}$ (Nelson and Setser~\cite{nelson90}, read from a search
snippet of their abstract) with \Krs's yield and B/C split, and the
electrons pump it
from \Krs{} at Kannari's $\wf{k.76.A}\,T_e^{\wf{k.76.B}}\exp(-\wf{k.76.C}/T_e)$
and return it by detailed balance. The two stepwise ionizations have no
measurement (Vriens~\cite{vriens65} is a calculation) and are bands at 0, up to
$\wf{k.stepS.hi}$ and $\wf{k.stepP.hi}$ on Kannari's forms;
\texttt{--ladder kannari} sets every rung band to Kannari's value.

\textbf{The cycle in the kinetic equation.} Kannari's
$A\,T_e^{1/2}\exp(-C/T_e)$ is exactly the Maxwellian rate of
$\sigma(\varepsilon)=\sigma_0(1-C/\varepsilon)$ above $C$, with
$\sigma_0=\wf{pop.sigma0}\unit{cm^2}$, so the pair is a population process
on \Krs{} with that pseudo cross section (threshold the ledger's gap,
\wm{gap5p}\unit{eV}) and its Klein-Rosseland reverse on \Krpp{}. On
\wf{pop.cells} cells to \wf{pop.emax}\unit{eV} it gives the printed law to
\wf{pop.k} and the detailed-balance ratio to \wf{pop.db}. With the pair as
LIGHT-side rates alone the kinetic energy balance missed LIGHT's by 0.7 of
the input; in a \wf{pop.short.t}-ns pulse with it inside, the field
carries \wf{pop.short.ratio.lo} to \wf{pop.short.ratio.hi} of the input and
the balances agree to \wf{pop.short.gap}.

\textbf{KrF(B,C) + F} (QF, QFC): a band, 0 to Kannari's `analogous to
F$_2$'~\cite{kannari85} $\wf{k.krff.hi}\cm{3}$. No measurement of it, or of any
RgX$^*$ + X analogue, exists.

\textbf{The optics.} $g_0=\sigma_B n_B$ with $\sigma_B=\wf{sig.B}\unit{cm^2}$
and $\alpha=\sum_i\sigma_i n_i$ over \wf{n.abs} absorbers: \Fm{}
$\wf{sig.Fm}$ (measured), \Krtp{} $\wf{sig.Kr2p}$ (an estimate), \Krpp{}
$\wf{sig.Kpp}$ (a calculation), \Ft{} $\wf{sig.F2}$, \Krs{} $\wf{sig.Krs}$
(compiled) and \KtF{} $\wf{sig.K2F}\unit{cm^2}$, Hakuta et al.'s direct
measurement~\cite{hakuta87} in place of the compilation's
$\wf{sig.K2F.comp}$~\cite{kannari85,johnson80}; the cross sections are
Kannari's Table I, \Fm{} from Mandl~\cite{mandl71}, \Krpp{} from
Hyman~\cite{hyman77}. A
diagnostic: LIGHT has no photon field, so the absorbers change no density.

\textbf{Ion-ion recombination.} Flannery's turnover (as Kannari~\cite{kannari85}
writes it) as the harmonic mean of REBUILD's three-body limb,
$\wf{ii.k3}\,N$ at 1 amagat, and the Langevin limb of the measured reduced
mobilities (Ellis et al.~\cite{ellis76,ellis78};
\Krp{} \wf{K0.Krp}, \Krtp{} \wf{K0.Kr2p}, \Fm{} \wf{K0.Fm}\unit{cm^2/V\,s}),
$\wf{ii.kL}/N$ for \Krp{} + \Fm{} and $\wf{ii.kLb}/N$ for \Krtp{} + \Fm{},
where REBUILD's was $\wf{ii.kL.v03}/N$ for both. The 1-amagat rate moves
from $\wf{ii.k1.v03}$ to $\wf{ii.k1}$ (\Krp) and $\wf{ii.k1b}$
(\Krtp)\cm{3}. The KrF$^*$ yield (ii\_yield, 1 by default, down to
\wf{ii.yield.lo}) and its C share (ii\_f\_c, 0 as in v0\_3, up to
\wf{ii.fc.hi}) are bands on the three product rows.

\textbf{The walls.} In a gap of \wf{wl.gap}\unit{cm} the cations leave at
the ambipolar rate of the electropositive global model~\cite{lieberman05}, $D_a/\Lam^2$, $D_a=\mu_i k(T_e+T_g)/e$ with the
mobilities above and \Lam{} the gap over $\pi$, \Fm{} at its free
diffusion; each lost pair takes its formation energy and \wf{wl.eTe}\,$T_e$
to the wall and \wf{wl.iTe}\,$T_e$ to the gas through the ambipolar
field. At the reference point's mid-pulse the rate is \wf{wl.nu}\unit{s^{-1}},
the walls take \wf{wl.share}\,\% of the input and the 248-nm share is
\wf{wl.248}\,\%; in a \wf{wl.thin.gap}-cm gap at \wf{wl.thin.p}\unit{atm}
they take \wf{wl.thin.share}\,\%. No 1-D number checks this closure yet:
1Dv0\_2's walls line at the same point is \wf{wl.1d}\,\% on
\wf{wl.1d.cells} uniform cells, but that line also holds its transport
step's mean-energy cut in the wall cells, and 1Dv0\_3, whose sheaths are
resolved, prints no walls line of its own.

\textbf{The wave.} A second input whose absorbed power is an output. The
electron-atom part is v0's Lorentz-gas kernel on the Biagi and SIGLO
momentum transfer~\cite{lxcat}, the electron-ion part Turnbull et al.'s
inverse bremsstrahlung~\cite{turnbull23,turnbull24} (the Sommerfeld
logarithm with the Gaunt-factor fit they take from Pradler and Semmelrock,
ion screening at the ion-sphere radius, the Langdon factor~\cite{langdon80}),
with their stated \wf{wv.err}\,\% error as the band ib\_scale
(\wf{ib.lo} to \wf{ib.hi}) on the electron-ion share of the power. Under
the kinetic closure the wave is Hagelaar and Pitchford's~\cite{hagelaar05}
high-frequency field term with Turnbull's per-energy cross section in its momentum
transfer, and for a Maxwellian handed to the solver that term gives the
kernel's atom part to \wf{hf.atoms} and its ion part to \wf{hf.ions} (at
$T_e=\wf{hf.te}\unit{eV}$, $n_e$ of \wf{hf.ne}\unit{cm^{-3}}). The ion
part here is the marginal one, the power with the ions in the momentum
transfer less the power without them, and that is what the band acts on
under both closures. It is not the kernel's printed electron-ion share,
$\nu_{ei}$'s share of the numerator: at 1\unit{atm} and 95\unit{GHz}
$\nu_{en}$ is \wf{hf.nuw} times $\omega$, the wave is collisional, and
the printed share is \wf{hf.share} times the marginal part (the two
agree only for $\nu_{en}\ll\omega$). Two double-counting rules: the
Langdon factor is applied under the Maxwellian closures only, and the
ions enter each field term once (the DC stand-in field sees them
Debye-screened, the wave with the IB logarithm). At the reference point a
wave of \wf{wv.E}\unit{V/cm} rms at \wf{wv.f}\unit{GHz} deposits
\wf{wv.em}\unit{J\,cm^{-3}} beside the \wf{wv.given} given
(\wf{wv.share}\,\% of all input), the band moving it by \wf{wv.band}\,\%;
the ions' marginal share at mid-pulse ($T_e$ \wf{wv.mid.te}\unit{eV},
$n_e$ \wf{wv.mid.ne}\unit{cm^{-3}}, $\nu_{en}/\omega$ \wf{wv.mid.nuw}) is
\wf{wv.mid.ei}\,\% (the printed share \wf{wv.mid.ei.share}\,\%) with a
Langdon factor of \wf{wv.mid.fL}, and at $10^{15}$ and
$10^{16}\unit{cm^{-3}}$ and 1.5\unit{eV} it is \wf{wv.ei.15.ghz} and
\wf{wv.ei.16.ghz}\,\% at 95\unit{GHz}, \wf{wv.ei.15.thz} and
\wf{wv.ei.16.thz}\,\% at 1\unit{THz}.
At 1\unit{THz} a field of \wf{wv.thz.E}\unit{kV/cm} absorbs
\wf{wv.thz.em}\unit{J\,cm^{-3}} in the microsecond. \texttt{ib.py} is held
to the IB project's own implementations by \wf{ib.checks} checks.

\section{How the numbers were made}

Every number in this note is written by
\texttt{docs/make\_overview\_v0\_4.py} from v0\_4 and from \texttt{ML/},
or by \texttt{docs/make\_system\_map\_v0\_4.py} (the reference point under
v0\_2's closure, its bands and its optics, which the system map draws), and
a test fails when those files no longer match the model. ML runs v0\_4
when \texttt{LIGHT\_VERSION} names it. Tahoma runs under the prototype beam
drive, a test fixture in \texttt{docs/} and not part of the model, in ML's
configuration (v0\_2's: the default network with every band at its default
and $n_e$ free from a seed of \wi{tah.seed}\unit{cm^{-3}}). The kinetic
EEDF takes minutes a run, so Tahoma's variants, its saturation intensity
and its formation efficiency run on v0\_2's closure here (the kinetic
closure moves v0\_4's Tahoma numbers by about 1\,\%, Table~\ref{tab:score}),
and the kinetic numbers are the reference point and Tahoma's default,
\wf{ref.min} and \wf{tah.kin.min} minutes each.

\section{Tahoma and the optics}

\begin{table}[H]
\centering\footnotesize
\setlength{\tabcolsep}{3.5pt}
\begin{tabular}{@{}llllll@{}}
\toprule
v0\_2's closure & $g_0$ (\%/cm) & $\alpha_0$ (\%/cm) & $g_0-\alpha_0$ & $T_e$ (eV) & $n_e$ (cm$^{-3}$)\\
\midrule
measured~\cite{tahoma89} & \wf{tah.g0.meas} & \wf{tah.alpha.meas} & \wf{tah.net.meas} & & \\
\wfTahVariants
v0\_3 & \wf{tah.v03.g0} (\x{\wf{tah.v03.g0.x}}) & \wf{tah.v03.alpha}
  (\x{\wf{tah.v03.alpha.x}}) & \wf{tah.v03.net} (\x{\wf{tah.v03.net.x}}) &
  \wf{tah.v03.te} & \wf{tah.v03.ne}\\
\midrule
v0\_4, kinetic EEDF & \wf{tah.kin.g0} (\x{\wf{tah.kin.g0.x}}) &
  \wf{tah.kin.alpha} (\x{\wf{tah.kin.alpha.x}}) & \wf{tah.kin.net}
  (\x{\wf{tah.kin.net.x}}) & \wf{tah.kin.te} & \wf{tah.kin.ne}\\
\bottomrule
\end{tabular}
\caption{Tahoma over the \wi{tah.win.lo} to \wi{tah.win.hi}\unit{ns}
window: v0\_4 with every band at its default and each variant alone, v0\_3,
and v0\_4 under the kinetic EEDF. The absorption by species with the
defaults: \wf{tah.by}\unit{\%/cm} (v0\_2's closure), \wf{tah.kin.by}
(kinetic); pure Kr, no \Ft: \wf{tah.pure.by}.}
\label{tab:tahoma}
\end{table}

The absorption is the first LIGHT has scored. With Kannari's compilation
value for \KtF{} it came out \wf{tah.comp.alpha}\unit{\%/cm},
\x{\wf{tah.comp.alpha.x}} the measurement, which sent the audit to the
literature: Hakuta, Komori, Mukai and Takuma~\cite{hakuta87} measured $\wf{sig.K2F}\unit{cm^2}$
at 248\unit{nm} by fluorescence suppression and wrote that it is `about 1/5
of the value assumed in the current kinetic model'. With the measurement
$\alpha_0$ is \x{\wf{tah.default.alpha.x}} of Tahoma's, \KtF{}
\wf{tah.k2f.share}\,\% of it, and the net gain \x{\wf{tah.default.net.x}}.
The pure-Kr loss Tahoma measured at 257.3\unit{nm}, \wf{tah.pure.meas}\unit{\%/cm},
LIGHT gives as \wf{tah.pure.alpha} with 248-nm cross sections, a rough
check (\Krtp's $\wf{sig.Kr2p}\unit{cm^2}$ is itself an estimate). Under the
beam the 5s $\leftrightarrow$ 5p cycle cools the electrons
(\wf{tah.default.te} against v0\_3's \wf{tah.v03.te}\unit{eV}) and lowers
$n_e$, and the gain rises by a few percent (Table~\ref{tab:tahoma}'s
factors are the model over the measurement: \x{\wf{tah.default.g0.x}}
with the cycle, \x{\wf{tah.nocycle.g0.x}} without). The rung's own bands
move it by under 1\,\% (\x{\wf{tah.kannari.g0.x}} at Kannari's values,
\x{\wf{tah.kpslo.g0.x}} with \Krpp{} + Kr at its low end); the bands that
move it are v0\_2's and the new quenches, the trimer's temperature law
(\x{\wf{tah.tri.g0.x}}), the ion-ion yield (\x{\wf{tah.iiyield.g0.x}}),
KrF(B,C) + F at Kannari's analogue (\x{\wf{tah.krff.g0.x}}), the ion-ion
C share (\x{\wf{tah.iifc.g0.x}}) and \Krts{} + F (\x{\wf{tah.kr2f.g0.x}}),
none of them with a measurement behind the band. In the window under the
kinetic EEDF LIGHT's
mean energy is the kinetic one at \wf{tah.kin.reached} update, held by a
field of at most \wf{tah.kin.en}\unit{Td}, and the two energy balances
agree to \wf{tah.kin.gap} of the beam's electron power. ML's formation run
gives $g_0$ \wf{tah.ml.g0}\unit{\%/cm} with its own channels, $I_\mathrm{sat}$
\wf{tah.isat}\unit{MW/cm^2} and \gIP{} $=$ \wf{tah.eta}.

\section{The reference point}

At the notes' reference point, \wm{p}\unit{atm}, \pc{xF2} \Ft{},
\wm{P}\unit{kW\,cm^{-3}} for \wm{pulse}\unit{\mu s} and run to
\wm{tend}\unit{\mu s}, v0\_4 gives these shares of the deposited energy:

\begin{table}[H]
\centering\footnotesize
\begin{tabular}{@{}lrrrr@{}}
\toprule
 & v0\_4, kinetic & v0\_4, v0\_2's closure & v0\_3, kinetic & v0\_3, v0\_2's closure\\
\midrule
248\unit{nm}, KrF B-X & \wf{ref.248} & \wm{out.248} & \wk{ref.248} & \wk{ref.248.v02}\\
272\unit{nm}, KrF B-A & \wf{ref.272} & \wm{out.272} & \wk{ref.272} & \wk{ref.272.v02}\\
275\unit{nm}, KrF C-A & \wf{ref.275} & \wm{out.275} & \wk{ref.275} & \wk{ref.275.v02}\\
420\unit{nm}, \KtF{} & \wf{ref.420} & \wm{out.420} & \wk{ref.420} & \wk{ref.420.v02}\\
147\unit{nm}, \Krts{} & \wf{ref.147} & \wm{out.147} & \wk{ref.147} & \wk{ref.147.v02}\\
near IR, 5p \ra 5s & \wf{ref.IR} & \wm{out.IR} & -- & --\\
all bands & \wf{ref.all} & \wm{notes.emit} & \wk{ref.all} & \wk{ref.all.v02}\\
heat, reactions & \wf{ref.heat} & \wm{out.heat} & & \\
to \Krs{}, gross & \wf{ref.krstar} & \wm{notes.krs} & \wk{ref.krstar} & \wk{ref.krstar.v02}\\
to \Krp{} & \wf{ref.krplus} & \wm{notes.krp} & \wk{ref.krplus} & \wk{ref.krplus.v02}\\
KrF(B) formed & \wf{ref.krfb} & \wm{krf.eff} & \wk{ref.krfb} & \wk{ref.krfb.v02}\\
elastic heating & \wf{ref.elastic} & \wm{out.elastic} & \wk{ref.elastic} & \wk{ref.elastic.v02}\\
$g_0$, 500 to 1000\unit{ns} (\%/cm) & \wf{ref.g0} & \wm{op.g0} & & \\
$\alpha$ (\%/cm) & \wf{ref.alpha} & \wm{op.alpha} & & \\
$T_g$ at the end (K) & \wf{ref.tg} & \wm{TgEnd} & \wk{ref.tg} & \wk{ref.tg.v02}\\
$T_e$, pulse mean (eV) & \wf{ref.te} & \wm{Te} & & \\
\bottomrule
\end{tabular}
\caption{The reference point, percent of \Edep{} unless a row says
otherwise. v0\_4's v0\_2-closure column is the system map's run.}
\end{table}

What moves the result is the 5s $\leftrightarrow$ 5p cycle. On Kannari's
fit the electrons pump \pc{brs.76} of the \Krs{} formed up to the rung
(\wf{ref.brs76}\,\% under the kinetic EEDF), the rung relaxes
\pc{brp.KPS} of it to \Krs{} on Kr and the \wm{gap5p}\unit{eV} of every
cycle heats the gas: the cycle takes \pc{cyc.e} of the input, cuts the
\Krs{} population and with it the \Krts{} route, and the 248-nm share rises
as the harpoon and the trimer share what is left. The near-IR 5p \ra 5s
band is \wm{out.IR}\,\%: the measured relaxation outruns the \wf{tau5p}-ns
radiation. With the rung neutralised (the cycle off, the rung relaxed at
once) the share is \wm{om.v03}\,\%, v0\_3's; with every rung band at
Kannari's value \wm{om.kannari}\,\%. Under the kinetic EEDF the field
carries \wf{ref.ratio.lo} to \wf{ref.ratio.hi} of the input through the
pulse's second half and the two energy balances agree to \wf{ref.gap};
the run takes \wf{ref.solutions} kinetic solutions, \wf{ref.on} in the
pulse, with $E/N$ between \wf{ref.en.lo} and \wf{ref.en.hi}\unit{Td}, and
closes its ledger below \wf{ref.closure}.

Every band's two ends under v0\_2's closure, one band at a time
(\texttt{run\_v0\_4.py --bands}):

\begin{table}[H]
\centering\scriptsize
\resizebox{\linewidth}{!}{%
\begin{tabular}{@{}lllllllr@{}}
\toprule
band & default & range & 248\unit{nm} & 147\unit{nm} & $g_0$ & $\alpha$ & $T_g$ (K)\\
\midrule
\wmBands
\bottomrule
\end{tabular}}
\caption{Each band alone at its two ends, the others at their defaults
(\wm{bd.n} bands, \wm{bd.nzero} with no measurement and a default of 0).
The columns are the 248-nm and 147-nm shares of \Edep{} (\%), the
500-to-1000-ns gain and absorption (\%/cm) and the final gas
temperature, low end to high end. With every band at its default the
248-nm share is \wm{out.248}\,\%; no band but the four named in the text
moves it by more than \wm{bd.minor} points.}
\label{tab:bands}
\end{table}

\section{The energy ladder}

v0\_4's ladder (Figure~\ref{fig:ladder}) adds the 5p rung, \lad{E.PP}\unit{eV}
above Kr and \lad{gap5p}\unit{eV} above \Krs{}, with K76 and K76r between
the two, KPS and K27X as quenches, PRAD's \lad{hv.PRAD}-eV photon, the
two stepwise ionizations and the harpoon from 5p; \Krtp{} + \el{} now
lands on it, leaving \lad{heat.R11p}\unit{eV} to the gas where v0\_3 left
2.93. KrF(B) lies \lad{DeB}\unit{eV} below the ion pair \Krp{} + \Fm{} and
\lad{BoverKrF}\unit{eV} above Kr + F, and the \lad{wl.G4}-nm photon
carries \lad{hv.G4}\unit{eV} of that.

\begin{figure}[p]
\centering
\ladpicture
\caption{The energy ladder of LIGHT v0\_4, version \lad{version}, in eV
above ground-state Kr and \Ft{} at rest. A bar is a species at its
formation energy and a dashed bar a pair, the sum of two. An arrow is a
row, labelled with its id. It is solid in the default network and dotted
for a band whose default is 0, \lad{zero}. The rows of the optional group
\lad{groups} are dashed. A photon's arrow drops by its quantum, and the red
squiggle below it is the heat its lower state gives the gas as it falls
apart. A red stub is a quench. The energy a row releases heats the gas,
except a photon's, which leaves, and that of \lad{toe}, which goes to the
electrons. KrF(B) and KrF(C) lie \lad{gapBC}\unit{cm^{-1}} apart and are
drawn as a pair of bars that every B and C row meets. The rows
\lad{notdrawn} move nothing between levels and are not drawn.}
\label{fig:ladder}
\end{figure}

\section{Where v0\_4 is not yet reliable}

\begin{itemize}
\item The 5s $\leftrightarrow$ 5p rate is a fit, not a measurement, and it
decides the partition. Jung et al.~\cite{jung05,jung06} measured the cross
sections out of the J = 2 metastable into the ten 5p levels; their
absolute values are not among the sources held. Band k\_5s5p runs to 0 (\wm{bd.k_5s5p.lo248} to
\wm{bd.k_5s5p.hi248}\,\% at 248\unit{nm}).
\item KPS is one level's total loss carried to the manifold, its product
unresolved; K27 is a search snippet of an abstract. Both are flagged in
the rows.
\item The stepwise ionizations and the 5p harpoon are LIGHT-side rates,
outside the Boltzmann equation (0 by default). \Krtp{} + \el{} makes the
rung entirely; Shiu and Biondi say `mostly'.
\item The absorbers are a diagnostic; \Krtp's and \Krpp's cross sections
are an estimate and a calculation, and the pure-Kr loss is checked with
248-nm values against a 257-nm measurement.
\item The wall loss is the electropositive global-model closure, stated and
reported. The wave's field is the one inside the plasma: no propagation,
reflection or skin depth.
\item Everything in v0\_3's list: the steady EEDF after the pulse, the
first cell, the thermal electrons' grid bias, \Krts{} + F, the trimer above
its measured range, Hakuta's and Burnham and Searles' conflicts, and the
pure-krypton VUV yield at low pump.
\end{itemize}

\appendix
\section{Reaction rates}

v0\_4's rows are v0\_2's (its overview's Appendix~A) and these, each
with the value the model ran at the reference point and its source, as
\texttt{v0\_4/network.py} cites it; \texttt{v0\_4/SOURCES.md} section 7
says how each was read.

\begin{table}[H]
\centering\scriptsize
\renewcommand{\arraystretch}{1.1}
\begin{tabularx}{\linewidth}{@{}>{\raggedright\arraybackslash}p{1.7cm}>{\raggedright\arraybackslash}p{3.9cm}>{\raggedright\arraybackslash}p{3.5cm}>{\raggedright\arraybackslash}X@{}}
\toprule
row & reaction & value & source\\
\midrule
C1a--c & \el{} + Kr \ra \el{} + \Krpp{} & Biagi groups~\cite{lxcat} & the higher excitations end on the rung (v0\_3: on \Krs)\\
R11p & \Krtp{} + \el{} \ra \Krpp{} + Kr & $\wm{k.R11p}\,(300\,\mathrm{K}/T_e)^{\wm{pow.R11p}}$ & Shiu and Biondi~\cite{shiu77}; the 5p lines in the afterglow\\
KPS & \Krpp{} + Kr \ra \Krs{} + Kr & $\wf{k.KPS}\cm{3}$ & Gazeli et al.~\cite{gazeli21}, 5p$'$[3/2]$_2$ in Kr; band to $\wf{k.KPS.lo}$ (Jancaitis~\cite{jancaitis84}, Ar analogue)\\
PRAD & \Krpp{} \ra \Krs{} + $h\nu$ & $\tau=\wf{tau5p}\unit{ns}$ & Dzier\.z\k{e}ga et al.~\cite{dzierzega00} Table 1, ten measured lifetimes\\
K76 & \el{} + \Krs{} \ra \el{} + \Krpp{} & $\wf{k.76.A}\,T_e^{\wf{k.76.B}}e^{-\wf{k.76.C}/T_e}$ & Kannari~\cite{kannari85} [76], Jacob and Mangano's fit~\cite{jacob76}; band to 0\\
K76r & \el{} + \Krpp{} \ra \el{} + \Krs{} & detailed balance, $g$ 12/36 & the statistical weights\\
KSTEP\_S & \el{} + \Krs{} \ra \Krp{} + 2\el{} & 0; band to $\wf{k.stepS.hi}$ & Kannari [79], Vriens's calculation~\cite{vriens65} ($\wf{k.stepS.kannari}$); REBUILD's template\\
KSTEP\_P & \el{} + \Krpp{} \ra \Krp{} + 2\el{} & 0; band to $\wf{k.stepP.hi}$ & Kannari [79], Vriens~\cite{vriens65}\\
K27, K27C, K27X & \Krpp{} + \Ft{} \ra KrF(B), KrF(C), Kr + 2F & $\wf{k.K27}\cm{3}$, split as \Krs's & Nelson and Setser~\cite{nelson90} (snippet); Kannari's 8.1e-10 a mark\\
QF, QFC & KrF(B,C) + F \ra Kr + 2F & 0; band to $\wf{k.krff.hi}$ & Kannari~\cite{kannari85}, `analogous to F$_2$'\\
K26, K108 & \Krp{}, \Krtp{} + \Fm{} \ra KrF(B) & $\wf{ii.k1}$, $\wf{ii.k1b}\cm{3}$ at 1 amagat & Flannery's turnover~\cite{kannari85}: REBUILD's three-body limb, the Langevin limb of Ellis's mobilities~\cite{ellis76,ellis78}\\
K26C, K108C, K26X, K108X & the C and dissociating shares & bands ii\_f\_c, ii\_yield & Rokni, Jacob and Mangano~\cite{rokni77} (relative); REBUILD's split\\
$\sigma$ at 248\unit{nm} & \Fm, \Krtp, \Krpp, \Ft, \Krs, \KtF & $\wf{sig.Fm}$, $\wf{sig.Kr2p}$, $\wf{sig.Kpp}$, $\wf{sig.F2}$, $\wf{sig.Krs}$, $\wf{sig.K2F}\unit{cm^2}$ & Kannari's Table I~\cite{kannari85} (its refs.\ 81, 82, 62, 84, 16: \cite{mandl71,rokni77,hyman77,johnson80}); Hakuta et al.~\cite{hakuta87} for \KtF\\
walls & $D_a/\Lam^2$ & $K_0$ \wf{K0.Krp}, \wf{K0.Kr2p}, \wf{K0.Fm}\unit{cm^2/V\,s} & Ellis et al.~\cite{ellis76,ellis78}\\
wave & Lorentz gas + Turnbull IB & band ib\_scale \wf{ib.lo} to \wf{ib.hi} & Turnbull et al.~\cite{turnbull23,turnbull24}; Hagelaar and Pitchford~\cite{hagelaar05}\\
\bottomrule
\end{tabularx}
\caption{The rows and constants v0\_4 adds. Rates in cm$^3$/s; a 0 is a
band whose default is 0, no measurement.}
\end{table}


\begin{thebibliography}{99}
\footnotesize
\setlength{\itemsep}{1pt}
\bibitem{eden78} Eden, Waynant, Searles and Burnham, Appl.\ Phys.\ Lett.\ 32, 733 (1978).
\bibitem{gadomski92} Gadomski, Xu, Setser and Nelson, Chem.\ Phys.\ Lett.\ 189, 153 (1992).
\bibitem{yu90} Yu and Setser, J.\ Phys.\ Chem.\ 94, 2934 (1990).
\bibitem{xu93} Xu, Gadomski and Setser, J.\ Chem.\ Phys.\ 99, 2591 (1993).
\bibitem{tahoma89} Spectra Technology, Inc., \emph{Experimental Support for KrF Laser Modeling}, final report to Los Alamos National Laboratory, 19 July 1989, OSTI 10136715 (Tables 3-2 and 3-3).
\bibitem{hakuta90} Hakuta, Fujino, Nakayama and Takuma, Jpn.\ J.\ Appl.\ Phys.\ 29, 1194 (1990).
\bibitem{burnham77} Burnham and Searles, J.\ Chem.\ Phys.\ 67, 5967 (1977).
\bibitem{morozov08} Morozov, Heindl, Kr\"ucken, Ulrich and Wieser, J.\ Appl.\ Phys.\ 103, 103301 (2008).
\bibitem{hakuta87} Hakuta, Komori, Mukai and Takuma, J.\ Appl.\ Phys.\ 61, 2113 (1987).
\bibitem{kannari85} Kannari, Obara and Fujioka, J.\ Appl.\ Phys.\ 57, 4309 (1985), Table I and Sec.\ II.B.
\bibitem{jacob76} Jacob and Mangano, Appl.\ Phys.\ Lett.\ 28, 724 (1976).
\bibitem{gazeli21} Gazeli et al., Phys.\ Plasmas 28, 043301 (2021).
\bibitem{dzierzega00} Dzier\.z\k{e}ga, Volz, Nave and Griesmann, Phys.\ Rev.\ A 62, 022505 (2000), Table 1.
\bibitem{nelson90} Nelson and Setser, Chem.\ Phys.\ Lett.\ 170, 430 (1990); read here as a search engine's rendering of the abstract.
\bibitem{shiu77} Shiu and Biondi, Phys.\ Rev.\ A 16, 1817 (1977).
\bibitem{vriens65} Vriens, Physica 31 (1965), as Kannari's ref.\ 79.
\bibitem{ellis76} Ellis, Pai, McDaniel, Mason and Viehland, At.\ Data Nucl.\ Data Tables 17, 177 (1976).
\bibitem{ellis78} Ellis et al., At.\ Data Nucl.\ Data Tables 22, 179 (1978).
\bibitem{rokni77} Rokni, Jacob and Mangano, Phys.\ Rev.\ A 16, 2216 (1977).
\bibitem{mandl71} Mandl, Phys.\ Rev.\ A 3, 251 (1971).
\bibitem{hyman77} Hyman, Appl.\ Phys.\ Lett.\ 31, 14 (1977).
\bibitem{johnson80} Johnson and Hunter, J.\ Appl.\ Phys.\ 51, 2406 (1980).
\bibitem{tellinghuisen76} Tellinghuisen, Hays, Hoffman and Tisone, J.\ Chem.\ Phys.\ 65, 4473 (1976).
\bibitem{grossl81} Gr\"ossl et al., J.\ Chem.\ Phys.\ 74, 1728 (1981).
\bibitem{jancaitis84} Jancaitis, \emph{Kinetic Modelling of Krypton Fluoride Laser Systems}, Lawrence Livermore National Laboratory report UCRL-53465 (1984), Table II.
\bibitem{jung05} Jung, Stone, Boffard, Anderson and Lin, Phys.\ Rev.\ Lett.\ 94, 163202 (2005).
\bibitem{jung06} Jung, Stone, Boffard, Anderson and Lin, Phys.\ Rev.\ A 73, 022722 (2006).
\bibitem{turnbull23} Turnbull et al., Phys.\ Rev.\ Lett.\ 130, 145103 (2023).
\bibitem{turnbull24} Turnbull et al., Phys.\ Plasmas 31, 063304 (2024).
\bibitem{langdon80} Langdon, Phys.\ Rev.\ Lett.\ 44, 575 (1980).
\bibitem{hagelaar05} Hagelaar and Pitchford, Plasma Sources Sci.\ Technol.\ 14, 722 (2005).
\bibitem{lxcat} The Biagi (Magboltz 7.1) and SIGLO databases, LXCat, www.lxcat.net.
\bibitem{lieberman05} Lieberman and Lichtenberg, \emph{Principles of Plasma Discharges and Materials Processing}, 2nd ed.\ (Wiley, 2005), Ch.\ 5.
\end{thebibliography}

\end{document}
